% main.tex - Informe Técnico TacticalGrid
% Compilar con: latexmk -lualatex -shell-escape main.tex

% preambulo.tex - Paquetes y configuraciones comunes para ambos documentos
\documentclass[12pt,a4paper,oneside]{report}

\usepackage{ucaldas}  % Nuestra plantilla personalizada
\usepackage{etoolbox}

% Bibliografía
\usepackage[backend=biber,style=ieee,sorting=none]{biblatex}
\addbibresource{bibliography.bib}

% Gráficos
\graphicspath{{images/}}

% Tablas compactas y ajustadas al ancho útil de A4.
\AtBeginEnvironment{tabular}{\small\sloppy\setlength{\tabcolsep}{3pt}}
\AtBeginEnvironment{longtable}{\small\sloppy\setlength{\tabcolsep}{3pt}\setlength{\LTleft}{0pt}\setlength{\LTright}{0pt}}

% Configuración de minted para Python
\setminted{
    fontsize=\small,
    breaklines=true,
    breakafter=,
    frame=lines,
    framesep=4pt,
    linenos=true,
    numbersep=8pt,
    tabsize=4,
    bgcolor=LightBlue
}

% Configuración de pgfgantt
\usepackage{pgfgantt}
\pgfplotsset{compat=1.18}

% Configuración pgf-pie
\usetikzlibrary{pie}

% Comandos personalizados para el proyecto
\newcommand{\tacticalgrid}{\textsc{TacticalGrid}}
\newcommand{\rf}[1]{RF-#1}
\newcommand{\rfnf}[1]{RNF-#1}
\newcommand{\sectionref}[1]{Sección~\ref{#1}}
\newcommand{\figureref}[1]{Figura~\ref{#1}}
\newcommand{\tableref}[1]{Tabla~\ref{#1}}
\newcommand{\listingref}[1]{Listing~\ref{#1}}

% Abreviaturas
\newcommand{\ie}{i.\,e.}
\newcommand{\eg}{e.\,g.}
\newcommand{\cf}{c.\,f.}
\newcommand{\etal}{\textit{et al.}}

% Glosario técnico
\newcommand{\g}[1]{\textit{#1}}
\newcommand{\code}[1]{\texttt{#1}}
\newcommand{\var}[1]{\textit{#1}}
\newcommand{\diagramfigure}[3]{%
    \begin{figure}[H]
    \centering
    \IfFileExists{images/#1.pdf}{%
        \includegraphics[width=#2\textwidth,height=0.72\textheight,keepaspectratio]{images/#1.pdf}%
    }{%
        \fbox{\parbox[c][4cm][c]{0.82\textwidth}{\centering Diagrama no disponible: \texttt{#1.pdf}}}%
    }%
    \caption{#3}%
    \end{figure}%
}

\newcommand{\screenshotfigure}[3]{%
    \begin{figure}[H]
    \centering
    \IfFileExists{images/screenshots/#1.png}{%
        \includegraphics[width=#2\textwidth]{images/screenshots/#1.png}%
    }{\IfFileExists{images/screenshots/#1.jpg}{%
        \includegraphics[width=#2\textwidth]{images/screenshots/#1.jpg}%
    }{%
        \fbox{\parbox[c][6cm][c]{0.82\textwidth}{\centering Captura pendiente: \texttt{#1.png / #1.jpg}\\[0.2cm]#3}}%
    }}%
    \caption{#3}%
    \end{figure}%
}

% Referencias a RF/RNF
\newcommand{\RF}[1]{RF-#1}
\newcommand{\RNF}[1]{RNF-#1}

% Comandos para algoritmos
\newcommand{\astar}{A\textsuperscript{*}}
\newcommand{\alphabeta}{$\alpha$-$\beta$}
\newcommand{\minimax}{Minimax}
\newcommand{\bigO}[1]{\mathcal{O}(#1)}

% Referencias a figuras/tablas
\newcommand{\Fig}[1]{Figura~\ref{#1}}
\newcommand{\Tab}[1]{Tabla~\ref{#1}}
\newcommand{\Lst}[1]{Listing~\ref{#1}}
\newcommand{\Eq}[1]{Ecuación~\ref{#1}}
\newcommand{\Sec}[1]{Sección~\ref{#1}}
\newcommand{\App}[1]{Apéndice~\ref{#1}}

% Colores para código
\definecolor{codebg}{HTML}{1A1B2E}
\definecolor{codefg}{HTML}{C4C4E0}
\definecolor{keyword}{HTML}{6C63FF}
\newcommand{\kw}[1]{\textcolor{keyword}{\textbf{#1}}}

% Para tablas anchas
\newcommand{\rotatetable}[1]{\begin{sidewaystable}#1\end{sidewaystable}}
\usepackage{rotating}

% Título para portada
\def\documenttitle{Informe Técnico}

\begin{document}

% Portada
% portada.tex - Portada común para ambos documentos
% Uso: % portada.tex - Portada común para ambos documentos
% Uso: % portada.tex - Portada común para ambos documentos
% Uso: \input{portada} en el documento principal

\begin{titlepage}
    \thispagestyle{empty}
    \centering
    \vspace*{1cm}
    
    % Logo UCALDAS
    \includegraphics[width=0.28\textwidth]{images/logo_ucaldas.png}%
    \par\vspace{0.8cm}
    
    % Universidad
    {\Large \textbf{UNIVERSIDAD DE CALDAS}\par}
    \vspace{0.3cm}
    {\large Facultad de Inteligencia Artificial e Ingenierías\par}
    {\large Departamento de Ingeniería de Sistemas y Computación\par}
    \vspace{0.6cm}
    
    % Línea decorativa
    {\color{ucaldasDorado}\rule{0.6\textwidth}{2pt}\par}
    \vspace{0.6cm}
    
    % Título del documento
    {\Huge \bfseries \color{ucaldasAzul} TacticalGrid\par}
    \vspace{0.5cm}
    \vspace{0.8cm}

    % Título según documento
    {\Large \bfseries \documenttitle\par}
    \vspace{0.5cm}
    {\large Sistemas Inteligentes I\par}
    \vspace{1cm}
    
    % Autores
    {\small
    \begin{tabular}{c}
        \textbf{Autores} \\
        Jorge Iván Garcia Torres \\
        Daniel Felipe Franco Rincón \\[0.15cm]
        \textbf{Docente} \\
        Jairo Iván Vélez Bedoya \\[0.15cm]
        \textbf{Materia} \\
        Sistemas Inteligentes I \\[0.15cm]
        \textbf{Entrega} \\
        Octubre de 2026 \\[0.15cm]
        \url{https://github.com/df1903/tactical-grid-ai}
    \end{tabular}\par}
    \vspace{0.8cm}
    
    % Logotipo adicional pequeño
    \vfill
    {\small \color{ucaldasGris}
    Manizales, Caldas -- Colombia \par
    2026\par}
\end{titlepage} en el documento principal

\begin{titlepage}
    \thispagestyle{empty}
    \centering
    \vspace*{1cm}
    
    % Logo UCALDAS
    \includegraphics[width=0.28\textwidth]{images/logo_ucaldas.png}%
    \par\vspace{0.8cm}
    
    % Universidad
    {\Large \textbf{UNIVERSIDAD DE CALDAS}\par}
    \vspace{0.3cm}
    {\large Facultad de Inteligencia Artificial e Ingenierías\par}
    {\large Departamento de Ingeniería de Sistemas y Computación\par}
    \vspace{0.6cm}
    
    % Línea decorativa
    {\color{ucaldasDorado}\rule{0.6\textwidth}{2pt}\par}
    \vspace{0.6cm}
    
    % Título del documento
    {\Huge \bfseries \color{ucaldasAzul} TacticalGrid\par}
    \vspace{0.5cm}
    \vspace{0.8cm}

    % Título según documento
    {\Large \bfseries \documenttitle\par}
    \vspace{0.5cm}
    {\large Sistemas Inteligentes I\par}
    \vspace{1cm}
    
    % Autores
    {\small
    \begin{tabular}{c}
        \textbf{Autores} \\
        Jorge Iván Garcia Torres \\
        Daniel Felipe Franco Rincón \\[0.15cm]
        \textbf{Docente} \\
        Jairo Iván Vélez Bedoya \\[0.15cm]
        \textbf{Materia} \\
        Sistemas Inteligentes I \\[0.15cm]
        \textbf{Entrega} \\
        Octubre de 2026 \\[0.15cm]
        \url{https://github.com/df1903/tactical-grid-ai}
    \end{tabular}\par}
    \vspace{0.8cm}
    
    % Logotipo adicional pequeño
    \vfill
    {\small \color{ucaldasGris}
    Manizales, Caldas -- Colombia \par
    2026\par}
\end{titlepage} en el documento principal

\begin{titlepage}
    \thispagestyle{empty}
    \centering
    \vspace*{1cm}
    
    % Logo UCALDAS
    \includegraphics[width=0.28\textwidth]{images/logo_ucaldas.png}%
    \par\vspace{0.8cm}
    
    % Universidad
    {\Large \textbf{UNIVERSIDAD DE CALDAS}\par}
    \vspace{0.3cm}
    {\large Facultad de Inteligencia Artificial e Ingenierías\par}
    {\large Departamento de Ingeniería de Sistemas y Computación\par}
    \vspace{0.6cm}
    
    % Línea decorativa
    {\color{ucaldasDorado}\rule{0.6\textwidth}{2pt}\par}
    \vspace{0.6cm}
    
    % Título del documento
    {\Huge \bfseries \color{ucaldasAzul} TacticalGrid\par}
    \vspace{0.5cm}
    \vspace{0.8cm}

    % Título según documento
    {\Large \bfseries \documenttitle\par}
    \vspace{0.5cm}
    {\large Sistemas Inteligentes I\par}
    \vspace{1cm}
    
    % Autores
    {\small
    \begin{tabular}{c}
        \textbf{Autores} \\
        Jorge Iván Garcia Torres \\
        Daniel Felipe Franco Rincón \\[0.15cm]
        \textbf{Docente} \\
        Jairo Iván Vélez Bedoya \\[0.15cm]
        \textbf{Materia} \\
        Sistemas Inteligentes I \\[0.15cm]
        \textbf{Entrega} \\
        Octubre de 2026 \\[0.15cm]
        \url{https://github.com/df1903/tactical-grid-ai}
    \end{tabular}\par}
    \vspace{0.8cm}
    
    % Logotipo adicional pequeño
    \vfill
    {\small \color{ucaldasGris}
    Manizales, Caldas -- Colombia \par
    2026\par}
\end{titlepage}
\clearpage

% Resumen
\chapter*{Resumen}
\addcontentsline{toc}{chapter}{Resumen}
% 00_resumen.tex - Resumen ejecutivo
% Incluido antes de la tabla de contenidos

\noindent \tacticalgrid\ es un entorno de experimentación para algoritmos de inteligencia artificial
construido sobre un juego táctico 2D por turnos. Dos bandos (A y B) compiten por un recurso
ubicado en zona neutral: una unidad debe recogerlo y llevarlo a su propia base antes que el
adversario. El juego no es un producto comercial; su propósito es servir de bancada de pruebas
para siete algoritmos fundamentales de IA: \textbf{BFS, DFS, UCS, \astar, Beam Search, Minimax y
Alpha-Beta}.

El proyecto implementa \textbf{todos los requisitos obligatorios} del enunciado: modelo de estado
compuesto inmutable, carga atómica de escenarios JSON con 9+ validaciones, siete algoritmos con
sus métricas (RF-15, RF-19, RF-22, RF-30), dos heurísticas admisibles (Manhattan escalada y
obstacle-aware BFS), función de utilidad con pesos configurables en JSON, verificación de
soluciones (6 checks RF-42), CLI con sweeps combinatorios y exportación CSV/JSON, GUI tkinter
con modo automático IA vs IA, y 10 casos de demostración con resultados esperado vs. obtenido.

\textbf{Estado:} Documento técnico final del proyecto. 281 tests pasan,
\texttt{ruff} + \texttt{mypy --strict} en verde. Código en \url{https://github.com/df1903/tactical-grid-ai}.
\clearpage

% Tabla de contenidos
\tableofcontents
\clearpage

% Lista de figuras y tablas
\listoffigures
\clearpage
\listoftables
\clearpage

% Glosario / Abreviaturas
\chapter*{Glosario y Abreviaturas}
\addcontentsline{toc}{chapter}{Glosario y Abreviaturas}
% glosario.tex - Glosario y Abreviaturas
% Incluido en main.tex antes del cuerpo principal

\begin{longtable}{p{3cm}p{3cm}p{8cm}}
\toprule
\textbf{Término} & \textbf{Abreviatura} & \textbf{Definición} \\
\midrule
\endfirsthead
\toprule
\textbf{Término} & \textbf{Abreviatura} & \textbf{Definición} \\
\midrule
\endhead
\midrule
\multicolumn{3}{r}{\textit{Continúa en la siguiente página}} \\
\endfoot
\bottomrule
\endlastfoot
Búsqueda en Anchura & BFS & Breadth-First Search: explora nivel por nivel, óptima en número de movimientos \\
Búsqueda en Profundidad & DFS & Depth-First Search: explora una rama hasta el final antes de retroceder \\
Costo Uniforme & UCS & Uniform-Cost Search: expande el nodo de menor costo acumulado $g(n)$ \\
A Estrella & \astar & A*: $f(n) = g(n) + h(n)$, busca óptimo con heurística admisible \\
Haz Local & Beam & Beam Search: mantiene solo los $k$ mejores candidatos por nivel \\
Minimax & MM & Algoritmo adversarial: MAX maximiza, MIN minimiza \\
Alfa-Beta & $\alpha$-$\beta$ & Minimax con poda: elimina ramas que no cambian la decisión \\
RF & RF & Requisito Funcional (del enunciado) \\
RNF & RNF & Requisito No Funcional \\
PLY & -- & Una acción de un bando (profundidad = número de plies) \\
MAX/MIN & -- & Roles en el árbol Minimax (no agentes distintos) \\
Utilidad & $U(s)$ & Valor numérico de un estado desde la perspectiva de MAX \\
Heurística Admisible & -- & $h(n) \le h^*(n)$: nunca sobreestima el costo real restante \\
Heurística Consistente & -- & $h(n) \le c(n,n') + h(n')$: cumple desigualdad triangular \\
Factor de Ramificación & $b$ & Número promedio de sucesores por estado \\
Profundidad de Solución & $d$ & Número de movimientos hasta la meta \\
Estado Compuesto & -- & Posiciones + carrier + resource + turn + turn\_number \\
\end{longtable}

\vspace{1cm}
\textbf{Abreviaturas de archivos:}
\begin{description}
\item[BFS] Breadth-First Search
\item[DFS] Depth-First Search
\item[UCS] Uniform Cost Search
\item[ASTAR] A* (A Estrella)
\item[BEAM] Beam Search
\item[MM] Minimax
\item[AB] Alpha-Beta
\item[CLI] Command Line Interface
\item[GUI] Graphical User Interface
\item[JSON] JavaScript Object Notation
\item[CSV] Comma-Separated Values
\item[SVG] Scalable Vector Graphics
\item[PDF] Portable Document Format
\item[UCALDAS] Universidad de Caldas
\item[IA] Inteligencia Artificial
\end{description}
\clearpage

% Cuerpo principal - numeración arábiga
\pagenumbering{arabic}

% CAPÍTULOS DEL INFORME TÉCNICO
% 01_introduccion.tex - Capítulo 1: Introducción
\chapter{Introducción}
\label{ch:introduccion}

\section{Contexto y Motivación}
\label{sec:contexto}

\tacticalgrid\ nace como proyecto integrador del curso \textbf{Sistemas Inteligentes I} (Universidad
de Caldas, 2026-2). El trabajo plantea la
construcción de un juego táctico 2D por turnos que sirva como \textit{entorno de experimentación}
para siete algoritmos clásicos de inteligencia artificial:

\begin{enumerate}
\item \textbf{BFS} (Breadth-First Search) -- búsqueda en anchura
\item \textbf{DFS} (Depth-First Search) -- búsqueda en profundidad
\item \textbf{UCS} (Uniform-Cost Search) -- costo uniforme
\item \textbf{A\textsuperscript{*}} -- A* con al menos dos heurísticas admisibles
\item \textbf{Beam Search} -- búsqueda con memoria limitada ($k = 1,2,4,8$)
\item \textbf{Minimax} -- decisión adversarial con profundidad limitada
\item \textbf{Alpha-Beta} -- Minimax con poda alfa-beta
\end{enumerate}

La premisa central del enunciado es clara: \textit{``El juego es el entorno de experimentación.
El proyecto es el sistema inteligente que decide dentro de él.''} No se busca un videojuego
comercial; la interfaz solo debe permitir \textbf{observar y experimentar con la inteligencia}
del sistema.

\section{Pregunta Central}
\label{sec:pregunta-central}

El enunciado plantea la pregunta guía del proyecto:

\begin{quote}
\textit{¿Cómo cambia la forma de decidir de un sistema inteligente cuando pasa de encontrar una
ruta, a optimizar su costo, utilizar conocimiento heurístico y finalmente actuar frente a un
adversario que también toma decisiones?}
\end{quote}

Esta pregunta estructura las cuatro etapas del proyecto:

\begin{enumerate}
\item \textbf{Encontrar una ruta} $\to$ BFS, DFS
\item \textbf{Optimizar el costo} $\to$ UCS
\item \textbf{Usar conocimiento heurístico} $\to$ A*, Beam Search
\item \textbf{Decidir frente a adversario} $\to$ Minimax, Alpha-Beta
\end{enumerate}

\section{Equipo y Entrega}
\label{sec:equipo}

\begin{table}[H]
\centering
\caption{Información del proyecto}
\begin{tabular}{ll}
\toprule
\textbf{Integrantes} & Jorge Iván Garcia Torres, Daniel Felipe Franco Rincón \\
\textbf{Materia} & Sistemas Inteligentes I -- Ingeniería de Sistemas y Computación \\
\textbf{Universidad} & Universidad de Caldas (UCALDAS) \\
\textbf{Entrega} & Octubre de 2026 \\
\textbf{Repositorio} & \url{https://github.com/df1903/tactical-grid-ai} \\
\textbf{Stack} & Python 3.14, \texttt{uv}, \texttt{pytest}, \texttt{ruff}, \texttt{mypy --strict} \\
\bottomrule
\end{tabular}
\end{table}

\section{Estructura del Documento}
\label{sec:estructura-doc}

Este informe técnico sigue la estructura del enunciado y cubre:

\begin{enumerate}
\item \textbf{Modelado y estado} (Cap. \ref{ch:modelado}) -- Estado compuesto, reglas, RF-10 a RF-13
\item \textbf{Reglas y costos} (Cap. \ref{ch:reglas}, \ref{ch:costos}) -- Reglas deterministas, RF-09
\item \textbf{Arquitectura} (Cap. \ref{ch:arquitectura}) -- 5 componentes desacoplados (RNF-05)
\item \textbf{Algoritmos de búsqueda} (Cap. \ref{ch:busqueda} -- \ref{ch:heuristicas}) -- BFS, DFS, UCS, A*, Beam
\item \textbf{Búsqueda adversarial} (Cap. \ref{ch:adversarial}) -- Minimax, $\alpha$-$\beta$, utilidad
\item \textbf{Experimentación} (Cap. \ref{ch:experimentos}) -- Metodología, casos 1-10, rendimiento
\item \textbf{Verificación y testing} (Cap. \ref{ch:verificacion}) -- 6 checks RF-42, 281 tests
\item \textbf{Formato JSON} (Cap. \ref{ch:json}) -- 9+ validaciones, campos opcionales
\item \textbf{Limitaciones} (Cap. \ref{ch:limitaciones}) -- 8 ítems documentados
\item \textbf{Uso de IA} (Cap. \ref{ch:ia}) -- Reporte de uso de IA generativa
\item \textbf{Reproducibilidad} (Cap. \ref{ch:reproducibilidad}) -- Instalación, tag Git, comandos
\end{enumerate}

Todos los resultados numéricos (casos 1-10, rendimiento) se derivaron \textbf{a mano antes de
ejecutar el código} y son verificados por tests automatizados (\texttt{tests/test\_cases\_f*.py}).
Los archivos CSV/JSON brutos están en \texttt{docs/results/}.
% 02_modelado_estado.tex - Capítulo 2: Modelado del Problema y Representación de Estados
\chapter{Modelado del Problema y Representación de Estados}
\label{ch:modelado}

Este capítulo describe el modelado del problema como espacio de estados (Unidad 2 del curso) y
la representación del estado compuesto que soporta tanto las búsquedas de un agente (F2, F3)
como la búsqueda adversarial (F4). Cumple con \rf{10}--\rf{13} y el caso 6.

\section{Principios de Diseño}
\label{sec:principios-diseno}

El modelado sigue tres principios fundamentales:

\begin{enumerate}
\item \textbf{Inmutabilidad:} \code{GameState} y \code{Board} son \code{@dataclass(frozen=True)}.
    La inmutabilidad garantiza que los estados puedan usarse como claves en \code{dict}/\code{set}
    (visited, frontier, parents) sin riesgo de mutación accidental.
\item \textbf{Separación estática/dinámica:} \code{Board} contiene todo lo que \textbf{nunca
    cambia} (grid, catálogo de terrenos, bases, bando de cada unidad, límite de turnos).
    \code{GameState} contiene solo lo que \textbf{cambia} (posiciones, carrier, resource, turn,
    turn\_number).
\item \textbf{Estado compuesto completo (RF-10, RF-11):} Un estado no es solo una posición;
    dos estados con la misma posición física son distintos si difieren en \code{carrier},
    \code{resource}, \code{turn}, \code{turn\_number} o las posiciones de otras unidades.
    El     \textbf{caso 6} demuestra esto explícitamente (ver \tableref{tab:caso6}).
\end{enumerate}

\section{Partes Estáticas: \code{Board}}
\label{sec:board}

\code{Board} es la parte estática del escenario; se construye una sola vez al cargar el JSON y
nunca muta. Contiene:

\begin{itemize}
\item \code{rows}, \code{cols}: dimensiones del grid
\item \code{grid}: tupla de tuplas con nombres de terreno por celda
\item \code{terrains}: \code{Mapping[str, Terrain]} -- catálogo \code{nombre $\to$ Terrain(cost, passable)}
\item \code{bases}: \code{Mapping[Side, Position]} -- base de cada bando
\item \code{unit\_sides}: \code{Mapping[str, Side]} -- bando de cada unidad (clave: id)
\item \code{turn\_limit}: límite de turnos para empate (default 100, configurable en JSON)
\end{itemize}

Métodos clave:
\begin{minted}{python}
def in_bounds(self, pos: Position) -> bool
def terrain_at(self, pos: Position) -> Terrain
def is_passable(self, pos: Position) -> bool
def cost_at(self, pos: Position) -> float
\end{minted}

\section{Parte Dinámica: \code{GameState}}
\label{sec:gamestate}

\code{GameState} es el estado compuesto inmutable (frozen dataclass). Contiene:

\begin{minted}{python}
@dataclass(frozen=True)
class GameState:
    positions: tuple[tuple[str, Position], ...]  # ordenado por unit_id
    carrier: str | None          # id de la unidad que lleva el recurso
    resource: Position | None    # celda del recurso si nadie lo lleva
    turn: Side                   # Side.A o Side.B
    turn_number: int = 0         # para límite de turnos / empate
\end{minted}

Métodos auxiliares:
\begin{minted}{python}
def position_of(self, unit_id: str) -> Position
def unit_at(self, pos: Position) -> str | None
\end{minted}

\section{Posición y Vecindad}
\label{sec:position}

\code{Position} es una tupla inmutable \code{(row, col)} con métodos:

\begin{minted}{python}
def neighbors(self) -> tuple[Position, ...]:  # 4 dirs: arriba, abajo, izq, der
def is_adjacent(self, other: Position) -> bool  # distancia Manhattan == 1
\end{minted}

La vecindad es de \textbf{4 direcciones} (arriba, abajo, izquierda, derecha) por decisión
del equipo (decisión 2026-10-07). Con 4 direcciones, la heurística Manhattan es admisible.

\section{Caso 6: Estado Compuesto Demostrado}
\label{sec:caso6}

El \textbf{caso 6} demuestra \rf{11}: ``dos estados con una unidad en la misma casilla pueden
ser distintos''. En el mapa del caso 7 (5$\times$7, coste 1), la unidad A1 en (1,0) junto a su
base (0,0):

\begin{table}[H]
\centering
\caption{Caso 6: mismo posición, estados distintos}
\begin{tabular}{lllll}
\toprule
\textbf{Estado} & \textbf{Pos A1} & \textbf{Carrier} & \textbf{Turn} & \textbf{Efecto de Move A1 $\to$ (0,0)} \\
\midrule
Libre (inicial) & (1,0) & nadie & A & entra a la base, no hay ganador \\
Portando & (1,0) & A1 & A & entra a la base $\to$ \textbf{gana A} \\
Turno rival & (1,0) & nadie & B & A1 no tiene acciones legales \\
\bottomrule
\end{tabular}
\end{table}
\label{tab:caso6}

Los tres estados tienen \code{positions} idénticas pero son objetos distintos
(\code{len(set(states)) == 3}). Un diccionario indexado por estado mantiene tres entradas;
indexado por posición solo mantendría una. Test: \texttt{tests/test\_case\_6.py}.

\section{Identificación Unívoca (RF-12)}
\label{sec:identificacion-univoca}

La Tabla~\ref{tab:identificacion} resume dónde se encuentra cada elemento requerido por \rf{12}.

\begin{table}[H]
\centering
\caption{Identificación unívoca de elementos del problema (RF-12)}
\begin{tabular}{ll}
\toprule
\textbf{Elemento} & \textbf{Dónde se define} \\
\midrule
Estado inicial & \code{Scenario.initial\_state} (construido por \code{scenario/loader.py}) \\
Acciones válidas & \code{domain.rules.legal\_actions(board, state)} \\
Estados sucesores & \code{domain.rules.successors(board, state)} $\to$ \code{(action, next\_state, cost)} \\
Condición de meta & Búsqueda: \code{problem.is\_goal(state)}; Juego: \code{winner(board, state)} \\
Costo de acción & \code{domain.rules.apply(board, state, action)} $\to$ \code{(next\_state, cost)} \\
\bottomrule
\end{tabular}
\label{tab:identificacion}
\end{table}

\section{Justificación de Suficiencia (RF-13)}
\label{sec:justificacion-suficiencia}

El estado contiene \textbf{exactamente} la información necesaria para decidir:

\begin{enumerate}
\item \textbf{Dónde está cada unidad} (\code{positions}) $\to$ legalidad de movimientos, ocupación
\item \textbf{De qué bando es cada unidad} (\code{unit\_sides} en \code{Board}) $\to$ turno, legalidad
\item \textbf{Quién lleva el recurso} (\code{carrier}) $\to$ pickup, delivery, interceptación
\item \textbf{Dónde está el recurso si nadie lo lleva} (\code{resource}) $\to$ pickup
\item \textbf{De quién es el turno} (\code{turn}) $\to$ qué unidades pueden moverse
\item \textbf{Número de turno} (\code{turn\_number}) $\to$ empate por límite
\end{enumerate}

Nada más influye en las reglas (no hay vida, munición, información oculta). Por tanto, dos
estados con iguales valores en estos seis campos tienen idéntico conjunto de acciones legales,
mismos sucesores, mismos costos y mismo resultado final. \rf{13} satisfecho.
% 03_reglas_juego.tex - Capítulo 3: Reglas del Juego
\chapter{Reglas del Juego}
\label{ch:reglas}

Este capítulo describe las reglas deterministas del juego (RF-09), decididas por el equipo
y documentadas en \texttt{decisions.md} (2026-10-07, 2026-10-14, 2026-10-18).

\section{Movimiento y Costos}
\label{sec:movimiento-costos}

\begin{itemize}
\item \textbf{Vecindad:} 4 direcciones (arriba, abajo, izquierda, derecha). Decisión 2026-10-07.
\item \textbf{Costo:} Se paga al \textbf{entrar} a la celda destino; la celda inicial no cuesta.
  Costos y transitabilidad provienen del catálogo \code{tipos\_terreno} del JSON.
\item \textbf{Wait:} Siempre legal, cuesta 0. No recoge recurso.
\end{itemize}

\section{Ocupación y Bases}
\label{sec:ocupacion-bases}

\begin{itemize}
\item \textbf{Una unidad por celda:} Validado al cargar (extra a RF-35).
\item \textbf{Base propia:} Se puede pisar; al llegar con el recurso $\to$ victoria inmediata.
\item \textbf{Base rival:} \textbf{NO} se puede pisar (decisión 2026-10-18).
\item \textbf{Celda recurso:} Se pisa; recoge al \textbf{llegar} (Move), no al esperar.
\end{itemize}

\section{Recurso e Interceptación (Placaje)}
\label{sec:recurso-interceptacion}

\begin{itemize}
\item \textbf{Pickup:} Al \textbf{llegar} a la celda del recurso (Move, no Wait).
\item \textbf{Entrega:} Llegar a base propia con recurso $\to$ victoria inmediata (salta interceptación).
\item \textbf{Interceptación (Placaje):} Tras \textbf{cualquier} movimiento, si un rival queda adyacente
  (distancia Manhattan = 1) al portador $\to$ recurso \textbf{cae en la celda del portador}
  (sin transferencia al rival; no hay desempate).
\item Portador original \textbf{no recoge} mientras siga en esa celda; debe salir y volver.
\item Cualquier \textbf{otra unidad} (incluso del mismo bando) lo recoge al llegar a la celda.
\item Se evalúa tras \textbf{cualquier} movimiento (de cualquier bando).
\end{itemize}

\section{Fin de Partida}
\label{sec:fin-partida}

\begin{itemize}
\item \textbf{Victoria:} Llevar recurso a base propia.
\item \textbf{Empate:} Límite de turnos (\code{juego.limite\_turnos}, default 100, configurable en JSON).
\item Bando sin movimientos legales $\to$ pasa turno (\code{Wait} siempre legal).
\end{itemize}

\section{Unidades}
\label{sec:unidades}

\begin{itemize}
\item Un solo tipo (\code{estandar}); mínimo 3 por bando en escenarios propios.
\item Apéndice A trae 2 $\to$ cargador no exige mínimo al cargar.
\end{itemize}

\section{Búsqueda de Un Agente (RF-14)}
\label{sec:busqueda-unagente}

Para búsquedas (BFS, DFS, UCS, A*, Beam):
\begin{itemize}
\item Rivales \textbf{removidos} del estado (si no, Apéndice A sin solución: B1 sobre meta).
\item Solo unidad inicio se mueve; compañeros quietos como obstáculos.
\item Turno fijo al bando de la unidad; recurso ignorado; \code{Wait} no usado.
\item Base rival sigue no entrable (regla de dominio).
\item Estado de búsqueda $\equiv$ posición de la unidad.
\end{itemize}

\section{Orden de Sucesores y Desempates}
\label{sec:orden-desempates}

\begin{itemize}
\item Orden fijo: \textbf{arriba, abajo, izquierda, derecha} (método \code{Position.neighbors()}).
\item DFS: hijos en orden \textbf{inverso} para expandir primero ``arriba''.
\item UCS/A*: desempate por \textbf{orden de entrada a la frontera} (FIFO tie-break).
\item Actualización de costo: camino más barato actualiza $g$ y padre, \textbf{no} regenera.
\item $max\_frontier$ = tamaño del set \code{waiting} (estados distintos esperando).
\end{itemize}
% 04_costos.tex - Capítulo 4: Costos
\chapter{Costos y Función de Costo}
\label{ch:costos}

\section{Definición de Costos (RF-03, RF-04)}
\label{sec:definicion-costos}

Los costos \textbf{nunca están hardcodeados} en el código (\rf{03}, \rf{37}, \rf{04}). Provienen
exclusivamente del catálogo \code{tipos\_terreno} del JSON cargado:

\begin{itemize}
\item \code{tipos\_terreno} es un mapeo \code{nombre $\to$ \{costo, transitable\}}
\item Solo terrenos \code{transitable: true} tienen \code{costo} (positivo, \code{float})
\item Terrenos no transitables (\code{transitable: false}) tienen \code{costo: null}
\item Acepta enteros o decimales (G12)
\end{itemize}

\section{Semántica del Costo}
\label{sec:semantica-costo}

\begin{itemize}
\item El costo se paga \textbf{al entrar} a la celda destino (decisión 2026-10-07).
\item La celda inicial \textbf{no cuesta}.
\item \code{Wait} cuesta \code{0.0} (decisión 2026-10-14).
\item No hay otras acciones con costo (solo \code{Move} y \code{Wait}).
\end{itemize}

Esto implica que $g(n)$ es la suma de los costos de las celdas \textbf{destino} de cada movimiento
en el camino. La celda inicial nunca contribuye a $g(n)$.

\section{Escala de Heurísticas (Sección 18.4 del Enunciado)}
\label{sec:escala-heuristicas}

Dado que los costos vienen del JSON y un mapa externo podría traer costos $< 1$, la heurística
Manhattan ``pura'' (1 por paso) dejaría de ser admisible si $\min(\text{costos}) < 1$.

La solución adoptada: \textbf{escalar por el costo transitable mínimo del escenario cargado}:

\begin{minted}{python}
def min_step_cost(board: Board) -> float:
    costs = [t.cost for t in board.terrains.values() if t.passable and t.cost is not None]
    return min(costs)
\end{minted}

Ambas heurísticas ($h_1$ Manhattan, $h_2$ obstacle-aware) multiplican su estimación en
pasos por \code{min\_step\_cost()}. Esto garantiza admisibilidad para \textbf{cualquier}
JSON válido, como se muestra en la definición anterior.

\section{Caso 2: Ruta Corta vs. Ruta Económica (RF-04)}
\label{sec:caso2}

El enunciado exige que existan mapas donde una ruta con menos movimientos cueste más que una
ruta más larga. \textbf{Caso 2} (\code{caso\_02\_bfs\_vs\_ucs.json}) lo demuestra:

\begin{itemize}
\item Banda de pantano (costo 7) bloquea el camino directo
\item BFS cruza el pantano: 7 movimientos, costo $4\times1 + 3\times7 = 25$
\item UCS rodea por terreno barato: 17 movimientos, costo 17
\end{itemize}

\begin{table}[H]
\centering
\caption{Caso 2: BFS vs UCS -- ruta corta vs. ruta económica}
\begin{tabular}{lccccc}
\toprule
Algoritmo & Costo & Movimientos & Generados & Expandidos & Max Frontera \\
\midrule
BFS & 25 & 7 & 126 & 104 & 23 \\
UCS & 17 & 17 & 269 & 239 & 33 \\
\bottomrule
\end{tabular}
\end{table}
\label{tab:caso2}


\section{Costo en Búsqueda de Un Agente (F2)}
\label{sec:costo-unagente}

En la abstracción de un agente (RF-14), el problema de búsqueda usa \code{PathProblem} donde
\code{expand} devuelve \code{[(GameState, cost)]} llamando a \code{domain.rules.successors} y
filtrando solo \code{Move} de la unidad inicio. El costo devuelto es \code{board.cost\_at(dest)}.

\section{Resumen de Reglas de Costo}
\begin{table}[H]
\centering
\caption{Resumen de reglas de costo}
\begin{tabular}{ll}
\toprule
\textbf{Regla} & \textbf{Valor} \\
\midrule
Costo de Move & \code{board.cost\_at(dest)} (terreno destino) \\
Costo de Wait & 0.0 \\
Celda inicial & No cuesta \\
Costos vienen de & \code{tipos\_terreno} del JSON (RF-03) \\
No hay costos hardcodeados & (RNF-03, RNF-04) \\
\bottomrule
\end{tabular}
\end{table}
\label{tab:resumen-costos}

% 05_arquitectura.tex - Capítulo 5: Arquitectura
\chapter{Arquitectura del Sistema}
\label{ch:arquitectura}

\section{Visión General}
\label{sec:arquitectura-vision}

La arquitectura sigue el principio de \textbf{separación de responsabilidades} (RNF-05): cinco
componentes desacoplados que se comunican mediante interfaces bien definidas. Ningún componente
conoce detalles de implementación de otro; solo dependen de interfaces (clases, funciones,
tipos de datos).

\begin{figure}[H]
\centering
\shorthandoff{>}
\begin{tikzpicture}[node distance=1.8cm, >=Stealth]
    \tikzset{
        component/.style={
            rectangle, draw=ucaldasAzul, thick, fill=ucaldasAzul!10,
            rounded corners, minimum width=2.5cm, minimum height=1.2cm,
            align=center, font=\small
        }
    }

    % Nodos
    \node[component] (JSON) {scenarios/*.json};
    \node[component, right=of JSON] (SCENARIO) {scenario/\\loader.py};
    \node[component, below=of SCENARIO] (DOMAIN) {domain/\\model.py, rules.py};
    \node[component, right=of DOMAIN] (ALGO) {algorithms/\\uninformed, informed,\\ adversarial, utility, heuristics};
    \node[component, below=of ALGO] (EXP) {experiments/\\runner, verifier, sweep};
    \node[component, right=of ALGO] (GUI) {gui/\\app.py, session.py};
    
    % Flechas
    \draw[->, thick] (JSON) -- node[above] {carga y valida} (SCENARIO);
    \draw[->, thick] (SCENARIO) -- node[left] {estado inicial + catálogo} (DOMAIN);
    \draw[->, thick] (DOMAIN) -- node[above] {sucesores, costo, meta} (ALGO);
    \draw[->, thick] (ALGO) -- node[right] {resultado + métricas} (EXP);
    \draw[->, thick, dashed] (EXP) -- node[left] {verificación} (ALGO);
    \draw[->, thick] (GUI) -- node[above] {orquesta} (SCENARIO);
    \draw[->, thick] (GUI) -- node[above] {ejecuta} (ALGO);
    \draw[->, thick] (GUI) -- node[above] {muestra} (EXP);
    
\end{tikzpicture}
\shorthandon{>}
\caption{Arquitectura de componentes y flujo de datos}
\label{fig:arquitectura}
\end{figure}

\section{Componentes}
\label{sec:componentes}

\begin{table}[H]
\centering
\caption{Componentes del sistema y responsabilidades}
\begin{tabular}{llp{5.5cm}}
\toprule
\textbf{Componente} & \textbf{Archivos clave} & \textbf{Responsabilidad} \\
\midrule
\code{domain} & \code{model.py}, \code{rules.py} & Estado compuesto, acciones, sucesores, costos, reglas de juego, ganador/empate \\
\code{scenario} & \code{loader.py} & Carga JSON, validación atómica (9+1), normalización, \code{TestConfig} \\
\code{algorithms} & \code{common.py}, \code{uninformed.py}, \code{informed.py}, \code{adversarial.py}, \code{utility.py}, \code{heuristics.py} & BFS, DFS, UCS, A*, Beam, Minimax, $\alpha$-$\beta$, heurísticas, utilidad \\
\code{experiments} & \code{runner.py}, \code{verifier.py}, \code{sweep.py} & Runner (modos), sweeps, export CSV/JSON, verificador RF-42 \\
\code{gui} & \code{app.py}, \code{session.py} & Ventana tkinter, dibujo, interacción, modo auto IA vs IA \\
\bottomrule
\end{tabular}
\label{tab:componentes}
\end{table}

\section{Comunicación entre Componentes}
\label{sec:comunicacion}

\begin{table}[H]
\centering
\caption{Comunicación entre componentes}
\begin{tabular}{llp{5.5cm}}
\toprule
\textbf{De} & \textbf{Hacia} & \textbf{Qué se comunica} \\
\midrule
\code{scenario} & \code{domain} & Estado inicial \code{GameState} + \code{Board} (catálogo validado) \\
\code{algorithms} & \code{domain} & Consulta \code{legal\_actions}, \code{successors}, \code{apply}, \code{winner} \\
\code{experiments} & \code{algorithms} & Ejecuta algoritmos, recibe \code{SearchResult}/\code{AdversarialResult} \\
\code{gui} & \code{scenario}, \code{algorithms}, \code{experiments} & Orquesta y muestra; \textbf{no contiene lógica de juego} \\
\bottomrule
\end{tabular}
\label{tab:comunicacion}
\end{table}

Principios clave:
\begin{itemize}
\item \code{algorithms} \textbf{nunca lee JSON} ni conoce la GUI; solo usa \code{domain} y \code{Scenario}.
\item \code{gui} no contiene lógica de juego; llama a \code{domain} (movimientos), \code{algorithms} (decisiones) y \code{experiments} (búsquedas).
\item \code{scenario} valida \textbf{atómicamente}: todo o nada, errores explicativos, nada parcialmente construido.
\end{itemize}

\section{Estructura de Carpetas}
\label{sec:estructura-carpetas}

\begin{minted}{text}
src/tactical_grid/
├── domain/         ← estado, acciones, sucesores, costos, reglas
├── scenario/       ← carga JSON, validación atómica
├── algorithms/     ← BFS, DFS, UCS, A*, Beam, Minimax, α-β, heurísticas, utilidad
├── gui/            ← interfaz gráfica tkinter
└── experiments/    ← métricas, verificador, casos mínimos, runner, sweep

scenarios/          ← JSON de escenarios (apendice_a.json, caso_01..07, prueba_15x15, etc.)
tests/              ← 281 pruebas automatizadas
docs/               ← enunciado, informe técnico, manual, casos, resultados, pdf/
docs/latex/         ← este documento (LaTeX)
\end{minted}

\section{Decisiones Arquitectónicas Clave}
\label{sec:decisiones-arquitectonicas}

\begin{table}[H]
\centering
\caption{Decisiones arquitectónicas registradas en \texttt{decisions.md}}
\begin{tabular}{ll}
\toprule
\textbf{Fecha} & \textbf{Decisión} \\
\midrule
2026-10-07 & 5 componentes: domain, scenario, algorithms, gui, experiments \\
2026-10-07 & Separación estática/dinámica: Board (estático) + GameState (dinámico) \\
2026-10-07 & Inmutabilidad total: frozen dataclasses, tuplas ordenadas \\
2026-10-07 & JSON claves en español, código en inglés \\
2026-10-16 & GUI: tkinter (estándar, cero dependencias) \\
2026-10-17 & CLI: sweep combinatorio, export CSV/JSON \\
\bottomrule
\end{tabular}
\label{tab:decisiones-arq}
\end{table}
% 06_algoritmos_busqueda.tex - Capítulo 6: Algoritmos de Búsqueda (F2, F3)
\chapter{Algoritmos de Búsqueda de Un Agente}
\label{ch:busqueda}

Este capítulo cubre las búsquedas de un solo agente (RF-14 a RF-24): BFS, DFS, UCS, A*, Beam Search.
El modelo de abstracción (\code{PathProblem}) transforma el juego completo en un problema de
un agente (RF-14).

\section{Abstracción de Un Agente: \code{PathProblem}}
\label{sec:pathproblem}

La RF-14 exige ignorar al adversario. \code{make\_problem} transforma el \code{GameState}
completo en un problema de \textbf{un solo agente}:

\begin{enumerate}
\item \textbf{Quita rivales:} las unidades del bando contrario se eliminan de \code{positions}.
\item \textbf{Compañeros quietos:} las otras unidades propias se quedan quietas como obstáculos.
\item \textbf{Turno fijo:} el turno se fija al bando de la unidad inicio; no alterna.
\item \textbf{Recurso ignorado:} \code{carrier=None}, \code{resource=None}.
\item \textbf{Turn number = 0:} se reinicia el contador.
\item \textbf{Base rival no entrable:} se mantiene la regla de dominio (no se puede pisar).
\end{enumerate}

\begin{minted}{python}
def make_problem(board: Board, initial: GameState, unit_id: str, goal: Position) -> PathProblem:
    side = board.unit_sides[unit_id]
    start = replace(
        initial,
        positions=tuple((uid, pos) for uid, pos in initial.positions
                        if board.unit_sides[uid] is side),
        turn=side,
        turn_number=0,
        carrier=None,
        resource=None,
    )
    return PathProblem(board, start, unit_id, goal)
\end{minted}

El estado de búsqueda equivale a la \textbf{posición de la unidad inicio}; los compañeros
permanecen como obstáculos estáticos.

\section{Convenciones Comunes (G7, G8, Decisión 2026-10-14)}
\label{sec:convenciones-comunes}

Todos los algoritmos de búsqueda (BFS, DFS, UCS, A*, Beam) comparten estas convenciones:

\begin{table}[H]
\centering
\caption{Convenciones comunes de búsqueda (G7, G8, decisión 2026-10-14)}
\begin{tabular}{ll}
\toprule
\textbf{Convención} & \textbf{Valor} \\
\midrule
Tipo de búsqueda & Grafo (visited set); un estado nunca se genera dos veces \\
Meta & Probada al \textbf{expandir} (no al generar) \\
Estado inicial & Cuenta como \textbf{generado} (cuenta = 1) \\
Orden de sucesores & Arriba, Abajo, Izquierda, Derecha (orden de \code{Position.neighbors()}) \\
DFS & Empuja hijos en orden \textbf{inverso} para expandir primero "arriba" \\
Desempate UCS/A* & Orden de entrada a la frontera (FIFO tie-break) \\
Actualización de costo & UCS/A*: camino más barato actualiza $g$ y padre, \textbf{no} regenera \\
$max\_frontier$ & Tamaño del conjunto \code{waiting} (estados distintos esperando) \\
DFS & Sin límite de profundidad (visited set garantiza terminación) \\
\bottomrule
\end{tabular}
\end{table}
\label{tab:convenciones-busqueda}


\textbf{Justificación:} La meta al expandir (no al generar) es necesaria para que UCS/A* garanticen
optimalidad. El estado inicial cuenta como generado porque entra a la frontera. El desempate
FIFO asegura determinismo y reproduce los resultados esperados de los casos de prueba.

\section{BFS: Búsqueda en Anchura}
\label{sec:bfs}

\begin{minted}{python}
def bfs(problem: PathProblem) -> SearchResult:
    # FIFO queue (deque)
    # goal test al expandir
    # no duplicados: skip si en explored o en parents (frontera)
    # orden: up, down, left, right
\end{minted}

\begin{itemize}
\item \textbf{Frontera:} \code{deque} (FIFO)
\item \textbf{Complejidad:} $\mathcal{O}(V)$ tiempo y espacio sobre $V$ estados alcanzables
\item \textbf{Óptimo en:} número de movimientos (si todos los costos son iguales)
\end{itemize}

\section{DFS: Búsqueda en Profundidad}
\label{sec:dfs}

\begin{minted}{python}
def dfs(problem: PathProblem) -> SearchResult:
    # LIFO stack (list)
    # hijos en orden inverso -> pop() expande primero "arriba"
    # mismo bookkeeping que BFS (grafo, meta al expandir, no duplicados)
\end{minted}

\begin{itemize}
\item \textbf{Frontera:} \code{list} (LIFO)
\item \textbf{Hijos en orden inverso:} para que \code{pop()} saque primero "arriba"
\item \textbf{Complejidad:} $\mathcal{O}(V)$ tiempo y espacio
\item \textbf{No óptimo} en movimientos ni en costo
\end{itemize}

\section{UCS: Costo Uniforme}
\label{sec:ucs}

Implementa la convención exacta de \rf{19} (curso): prioridad = menor $g(n)$; desempate por
orden de entrada a la frontera; actualización de costo sin regenerar.

\begin{minted}{python}
def ucs(problem: PathProblem) -> SearchResult:
    heap = [(0.0, 0)]           # (g, order_rank)
    order = {start: 0}          # estado -> rank de entrada
    by_order = [start]          # rank -> estado
    waiting = {start}           # estados en frontera (para max_frontier)
    
    while heap:
        g_cost, rank = heapq.heappop(heap)
        node = by_order[rank]
        if node in explored or g_cost > g[node]: continue  # stale entry
        waiting.discard(node)
        # goal test, expandir, actualizar...
        
        if child not in g:
            # PRIMERA VEZ
            order[child] = len(by_order)
            by_order.append(child)
            waiting.add(child)
            generated += 1
            heapq.heappush(heap, (new_cost, order[child]))
        elif new_cost < g[child]:
            # CAMINO MÁS BARATO: actualiza g y padre, MANTIENE order original
            g[child] = new_cost
            heapq.heappush(heap, (new_cost, order[child]))  # nuevo entry, mismo rank
\end{minted}

Puntos clave de la implementación:
\begin{enumerate}
\item \code{order} = orden de primera entrada a la frontera (tie-break RF-19)
\item Actualización de costo $\to$ \code{heappush} nuevo entry con \textbf{mismo rank}; el stale se descarta al extraer
\item \code{max\_frontier} = \code{len(waiting)} (solo estados distintos esperando, no entries stale)
\item \code{generated} solo incrementa en \textbf{primera generación} del estado
\end{enumerate}

\section{Métricas Reportadas (RF-15, RF-16, RF-19, RF-38)}
\label{sec:metricas-busqueda}

Todos los algoritmos de búsqueda devuelven \code{SearchResult} con campos equivalentes al
Apéndice B del enunciado:

\begin{table}[H]
\centering
\caption{Métricas de búsqueda (Apéndice B + tiempo)}
\begin{tabular}{ll}
\toprule
\textbf{Campo} & \textbf{Descripción} \\
\midrule
\code{algorithm} & Nombre: \code{BFS}, \code{DFS}, \code{UCS}, \code{A\_STAR\_H1}, \code{A\_STAR\_H2}, \code{BEAM\_K1}... \\
\code{success} & \code{True}/\code{False} \\
\code{path} & Tupla de \code{Position} (inicio a meta, inicio incluido) \\
\code{cost} & Costo total $g(goal)$ \\
\code{length} & Número de movimientos (\code{len(path)-1}) \\
\code{generated} & Estados incorporados a la búsqueda \\
\code{expanded} & Estados cuyos sucesores fueron analizados \\
\code{max\_frontier} & Máximo tamaño de \code{waiting} \\
\code{time\_seconds} & \code{time.perf\_counter()}; tiempo de ejecución \\
\bottomrule
\end{tabular}
\end{table}
\label{tab:metricas-busqueda}

% 07_heuristicas.tex - Capítulo 7: Heurísticas (F3)
\chapter{Heurísticas para A* y Beam Search}
\label{ch:heuristicas}

\section{Requisitos (RF-21, RF-22)}
\label{sec:requisitos-heuristicas}

El enunciado exige \textbf{al menos dos heurísticas} para A*:

\begin{enumerate}
\item \textbf{Heurística 1 (h1):} Distancia Manhattan (obligatoria)
\item \textbf{Heurística 2 (h2):} Propuesta y justificada por el equipo (RF-21)
\end{enumerate}

Para cada heurística se debe justificar (RF-22):
\begin{enumerate}
\item Qué intenta estimar
\item Qué información del estado utiliza
\item Cuál es su costo computacional
\item Si se considera admisible y por qué
\item En qué escenarios orienta bien y en cuáles pierde capacidad de discriminación
\end{enumerate}

Ambas heurísticas se escalan por \code{min\_step\_cost()} (costo transitable mínimo del JSON)
para garantizar admisibilidad aunque el mapa externo traiga costos $< 1$ (enunciado §18.4).

\section{h1: Manhattan Escalada}
\label{sec:h1}

\begin{minted}{python}
def manhattan(problem: PathProblem) -> Heuristic:
    scale = min_step_cost(problem.board)
    goal = problem.goal
    def h(state: GameState) -> float:
        pos = state.position_of(problem.unit_id)
        return (abs(pos.row - goal.row) + abs(pos.col - goal.col)) * scale
    return h
\end{minted}

\begin{table}[H]
\centering
\caption{Justificación RF-22 para h1 (Manhattan)}
\begin{tabular}{lp{8cm}}
\toprule
\textbf{Punto RF-22} & \textbf{Respuesta} \\
\midrule
1. Qué estima & Costo restante si no hay obstáculos y cada paso cuesta el mínimo \\
2. Info del estado & Posición de la unidad, meta, \code{min\_step\_cost} \\
3. Costo computacional & O(1) por evaluación, sin preprocesamiento \\
4. Admisible & Sí. Cada movimiento cambia la distancia Manhattan en $\le 1$ y cuesta $\ge$ min\_cost \\
5. Dónde pierde & Detrás de barreras: apunta recto al objetivo aunque haya muro (caso 3) \\
\bottomrule
\end{tabular}
\end{table}
\label{tab:h1-rf22}


\section{h2: Obstacle-Aware (BFS Inverso desde la Meta)}
\label{sec:h2}

Propuesta del equipo (RF-21): una heurística que \textbf{conoce el layout del mapa} (muros y
base rival) pero ignora costos de terreno y otras unidades.

\begin{minted}{python}
def obstacle_aware(problem: PathProblem) -> Heuristic:
    board = problem.board
    scale = min_step_cost(board)
    rival_base = board.bases[board.unit_sides[problem.unit_id].opponent]
    moves: dict[Position, int] = {problem.goal: 0}
    queue: deque[Position] = deque([problem.goal])
    while queue:
        cell = queue.popleft()
        for nxt in cell.neighbors():
            if nxt in moves or nxt == rival_base or not board.is_passable(nxt):
                continue
            moves[nxt] = moves[cell] + 1
            queue.append(nxt)
    def h(state: GameState) -> float:
        steps = moves.get(state.position_of(problem.unit_id))
        return math.inf if steps is None else steps * scale
    return h
\end{minted}

Construcción: \textbf{BFS inverso desde la meta} (backward BFS) esquivando muros y base rival.
Complejidad: O(R$\times$C) una vez al construir; O(1) por evaluación (lookup en dict).

\begin{table}[H]
\centering
\caption{Justificación RF-22 para h2 (Obstacle-aware)}
\begin{tabular}{lp{8cm}}
\toprule
\textbf{Punto RF-22} & \textbf{Respuesta} \\
\midrule
1. Qué estima & Costo restante con layout real (muros + base rival), paso a costo mínimo \\
2. Info del estado & Posición de la unidad + layout completo del mapa (muros, base rival) \\
3. Costo computacional & O(R$\times$C) una vez (BFS inverso), O(1) por evaluación (lookup) \\
4. Admisible & Sí. Es el costo exacto de un problema relajado (movimientos gratis excepto min\_cost) \\
5. Dónde pierde & No ve costos de terreno: si el atajo cruza pantano, h2 sigue prefiriéndolo (caso 4) \\
\bottomrule
\end{tabular}
\end{table}
\label{tab:h2-rf22}


\section{Propiedades: h2 Domina a h1}
\label{sec:h2-domina}

\begin{itemize}
\item h2 $\ge$ h1 \textbf{siempre} (h2 cuenta movimientos reales esquivando obstáculos; h1 asume camino recto)
\item Ambas son \textbf{admisibles} y \textbf{consistentes} para movimientos de 4 direcciones
\item h2 tiene O(R$\times$C) preprocesamiento una vez; h1 es O(1) sin preprocesamiento
\item En A*, h2 expande \textbf{menos estados} que h1 (casos 3 y 4 lo demuestran)
\item Celdas inalcanzables $\to$ h2 = $\infty$ (admisible: no hay camino, costo real = $\infty$)
\end{itemize}

\section{Medición de Efecto (Casos 3 y 4)}
\label{sec:medicion-heuristicas}

\begin{table}[H]
\centering
\caption{Expandidos en Caso 3 (muro) y Caso 4 (pantano vs. camino abierto)}
\begin{tabular}{lccc}
\toprule
Algoritmo & Caso 3 (muro) & Caso 4 (pantano vs. abierto) & Nota \\
\midrule
UCS & 295 & 332 & baseline \\
A* h1 (Manhattan) & 166 & 186 & h1 ignora muro/pantano \\
A* h2 (Obstacle-aware) & 84 & 115 & h2 ve muro, no ve pantano \\
\bottomrule
\end{tabular}
\end{table}
\label{tab:heuristicas-expandidos}


En caso 3 (muro vertical), h2 reduce expandidos a la \textbf{mitad} de h1 (84 vs 166).
En caso 4 (atajo por pantano), h2 ayuda pero no ve el costo del terreno.

\section{Origen de h2}
\label{sec:origen-h2}

La heurística h2 fue \textbf{propuesta por Claude Code (Sonnet 5.5)} durante el desarrollo;
el equipo revisó la admisibilidad, la consistencia y la dominancia sobre h1, y la aceptó
tras verificar que cumple los 5 puntos de RF-22. (La documentación inicial la llamaba
``propuesta del equipo''; se corrigió en \texttt{decisions.md} 2026-10-17).
% 08_busqueda_adversarial.tex - Capítulo 8: Búsqueda Adversarial (F4)
\chapter{Búsqueda Adversarial: Minimax y Alfa-Beta}
\label{ch:adversarial}

Este capítulo cubre la transición de ``encontrar una ruta'' a ``decidir frente a un adversario''
(RF-25 a RF-33). Se implementan Minimax y alfa-beta con profundidad configurable, utilidad
propia y 7 métricas RF-30.

\section{Cambio de Paradigma: De Búsqueda a Decisión}
\label{sec:cambio-paradigma}

En búsquedas F2/F3 (un agente): ``¿Cuál es mi mejor camino?'' $\to$ \code{PathProblem},
adversario ignorado.

En búsqueda adversarial F4: ``¿Cuál es mi mejor acción considerando las respuestas del
adversario?'' $\to$ \code{GameState} completo, \textbf{ambos bandos razonan}.

\section{Profundidad y Ply (G11, Decisión 2026-10-17)}
\label{sec:profundidad-ply}

\begin{itemize}
\item \textbf{1 unidad de profundidad = 1 acción de 1 bando} (1 ply)
\item Los bandos alternan por ply: MAX (lado que mueve en la raíz) $\to$ MIN $\to$ MAX $\to$ MIN...
\item \code{Wait} siempre es legal y cuenta como ply (consume el turno)
\item Profundidad configurable vía JSON: \code{juego.profundidad\_maxima\_minimax} o CLI \code{--depth}
\item Profundidad mínima = 1
\end{itemize}

\section{Función de Utilidad (RF-27, RF-33)}
\label{sec:utilidad}

Definida en \code{algorithms/utility.py}. Perspectiva: \textbf{MAX = bando que mueve en la raíz};
MAX maximiza, MIN minimiza. Pesos leídos del JSON (\code{utilidad}), defaults 10/3/1/0.5.

\begin{minted}{python}
U = w_entrega  * (+1 si MAX entregó, -1 si MIN entregó)
  + w_carga    * (+1 si MAX porta, -1 si MIN porta)
  - w_dist     * (distancia al objetivo actual)
  + w_rival    * (distancia del rival más cercano al portador, SOLO si MAX porta)
\end{minted}

\textbf{Objetivo dinámico (distancia):}
\begin{itemize}
\item MAX porta $\to$ objetivo = base propia de MAX
\item MIN porta $\to$ objetivo = posición del portador (perseguir al rival)
\item Nadie porta $\to$ objetivo = celda del recurso
\end{itemize}
Distancia = Manhattan desde la unidad MAX más cercana (desde el portador si MAX porta).

\textbf{Término rival:} Solo cuando MAX porta. Un rival cercano al portador de MAX \textbf{baja}
$U$ (riesgo de interceptación). Si MIN porta, no se suma (la persecución ya cubre ese riesgo).

\textbf{Pesos (desde JSON \code{utilidad}, defaults 10/3/1/0.5):}
\begin{table}[H]
\centering
\caption{Pesos de utilidad (RNF-03: configurables en JSON)}
\begin{tabular}{lll}
\toprule
\textbf{Peso} & \textbf{Default} & \textbf{Significado} \\
\midrule
\code{entrega} & 10 & Victoria/derrota domina todo \\
\code{recurso\_cargado} & 3 & Porta recurso > estar cerca \\
\code{distancia\_objetivo} & 1 & Cada paso de Manhattan = -1 \\
\code{distancia\_rival} & 0.5 & Rival cerca = riesgo (solo si MAX porta) \\
\bottomrule
\end{tabular}
\end{table}
\label{tab:pesos-utilidad}


\textbf{Terminal gana:} \code{winner()} devuelve $\pm$ \code{w\_entrega} ($\pm$10) y corta la
recursión. Victoria/derrota domina sobre cualquier combinación de distancias.

\section{Minimax y Alfa-Beta: Recursión Compartida}
\label{sec:minimax-alfabeta}

Ambos algoritmos comparten \textbf{una sola función recursiva} \code{value\_of};
alfa-beta solo añade el corte \code{alpha >= beta}.

\begin{minted}{python}
def value_of(node, level, alpha, beta):
    if level == depth or is_terminal(board, node):
        evaluated += 1
        return utility(node)
    maximizing = node.turn is max_side
    best = -inf if maximizing else inf
    for action in ordered(node):  # orden configurable: default/reversed
        child, _ = apply(board, node, action)
        generated += 1
        score = value_of(child, level + 1, alpha, beta)
        if maximizing:
            best = max(best, score)
            alpha = max(alpha, best)
        else:
            best = min(best, score)
            beta = min(beta, best)
        if prune and alpha >= beta:
            pruned += len(actions) - index - 1  # hermanos omitidos
            break
    return best
\end{minted}

\section{Métricas RF-30 (Implementadas en \code{AdversarialResult})}
\label{sec:metricas-adversarial}

\begin{table}[H]
\centering
\caption{Métricas adversariales (RF-30)}
\begin{tabular}{ll}
\toprule
\textbf{Métrica} & \textbf{Qué registra} \\
\midrule
\code{generated} & Hijos producidos (alfa-beta genera perezosamente: solo hasta corte) \\
\code{evaluated} & Estados donde se calculó utilidad (hojas + terminales) \\
\code{pruned} & Hermanos omitidos por \code{alpha >= beta} \\
\code{depth\_reached} & Nivel máximo alcanzado en la recursión \\
\code{time\_seconds} & \code{time.perf\_counter()}; tiempo de ejecución \\
\code{action} & Acción elegida (primera en orden con score estrictamente mayor) \\
\code{value} & Valor Minimax de la raíz \\
\code{root\_values} & Tupla \code{(acción, valor)} por acción raíz; en $\alpha$-$\beta$ los posteriores a la mejor pueden ser \textbf{cotas} \\
\code{considered} & (si \code{record=True}) Tupla \code{(GameState, Action)} de cada acción considerada (check 6) \\
\bottomrule
\end{tabular}
\end{table}
\label{tab:metricas-adversarial}


\section{RF-31: Misma Decisión con Mismo Orden}
\label{sec:rf31}

En la raíz, empate en score $\to$ se mantiene la \textbf{primera acción en el orden dado}.
Tanto Minimax como $\alpha$-$\beta$ iteran las acciones raíz en el mismo orden
(\code{default} o \code{reversed}) y actualizan \code{chosen} solo si \code{score > best\_value}
(estrictamente mayor). Por tanto:

\begin{itemize}
\item Misma profundidad, misma utilidad, mismo orden $\to$ \textbf{misma decisión}
\item Valores de acciones posteriores a la mejor en $\alpha$-$\beta$ pueden ser \textbf{cotas}
  (no valores exactos), porque el corte impide explorar completamente
\end{itemize}

\section{RF-32: Orden de Acciones y Poda}
\label{sec:rf32}

Orden de acciones configurable: \code{default} (arriba/abajo/izq/der por unidad) o
\code{reversed}. Cambiar el orden:

\begin{itemize}
\item \textbf{No cambia la decisión} (RF-31) si la utilidad y profundidad son iguales
\item \textbf{Sí cambia \code{pruned} y \code{generated}} (RF-32)
\item En caso 9 (profundidad 3): \code{default} genera 142, \code{reversed} genera 191
\item En caso 9 (profundidad 4): orden \code{reversed} cambia el tie-break entre valores
  iguales (0), produciendo acción diferente con mismo valor
\end{itemize}

\section{Check 6 del Verificador (RF-42)}
\label{sec:check6}

\code{verify\_adversarial(board, state, result)} re-ejecuta con \code{record=True} y valida:

\begin{enumerate}
\item Cada acción en \code{result.considered} es legal según \code{legal\_actions}
\item La acción elegida (\code{result.action}) es legal
\item No se generan estados inválidos durante la exploración
\end{enumerate}

Esto completa los 6 checks de RF-42 (5 para búsquedas + 1 para adversarial).

\section{Casos 7--10: Resultados}
\label{sec:casos7-10}

Escenario: \code{caso\_07\_decision\_adversarial.json} (5$\times$7, coste 1, 2 unidades/bando,
recurso (0,1), A mueve primero, pesos 10/3/1/0.5).

\begin{table}[H]
\centering
\caption{Casos 7--10: Resultados Minimax vs Alfa-Beta}
\begin{tabular}{lccccccc}
\toprule
Alg & Prof & Orden & Acción & Valor & Gen & Eval & Podados \\
\midrule
MM & 1 & def & A2$\to$(0,1) & 3 & 7 & 7 & 0 \\
MM & 2 & def & A2$\to$(0,1) & 0 & 71 & 64 & 0 \\
MM & 3 & def & A1$\to$(1,1) & 3 & 538 & 467 & 0 \\
MM & 4 & def & A1$\to$(0,0) & 0 & 4,826 & 4,296 & 0 \\
AB & 1 & def & A2$\to$(0,1) & 3 & 7 & 7 & 0 \\
AB & 2 & def & A2$\to$(0,1) & 0 & 31 & 24 & 40 \\
AB & 3 & def & A1$\to$(1,1) & 3 & 142 & 110 & 106 \\
AB & 4 & def & A1$\to$(0,0) & 0 & 431 & 299 & 757 \\
AB & 3 & rev & A1$\to$(1,1) & 3 & 191 & 148 & 145 \\
\bottomrule
\end{tabular}
\end{table}
\label{tab:casos7-10}


\begin{itemize}
\item \textbf{Caso 7:} La mejor decisión no es ir directo al recurso (profundidad 1 y 2 agarran;
  profundidad 3 protege moviendo A1$\to$(1,1) para bloquear a B1).
\item \textbf{Caso 8:} En profundidad 4, $\alpha$-$\beta$ genera \textbf{11$\times$ menos nodos}
  que Minimax (431 vs 4,826).
\item \textbf{Caso 9:} Mismo valor (3), misma decisión, pero orden \code{reversed} poda más
  (191 vs 142 generados en d=3). En d=4, orden \code{reversed} cambia tie-break (acción
  diferente, mismo valor 0).
\item \textbf{Caso 10:} Profundidades impares terminan en turno de A (valor 3, guarda);
  pares terminan en turno de B (valor 0, aplana). Efecto horizonte documentado en \S\ref{sec:limitaciones-conocidas}.
\end{itemize}
% 09_funcion_utilidad.tex - Capítulo 9: Función de Utilidad (detalle)
\chapter{Función de Utilidad: Detalle y Análisis}
\label{ch:utilidad}

Este capítulo profundiza en la función de utilidad $U(s)$, su diseño, pesos, limitaciones y
cómo explicarla en sustentación.

\section{Formulación Completa}
\label{sec:formulacion-completa}

La utilidad siempre se mide desde la perspectiva de \textbf{MAX = bando que mueve en la raíz}.

\begin{equation}
U(s) = w_{ent} \cdot D_{ent}(s) + w_{car} \cdot C_{car}(s) - w_{dist} \cdot D_{goal}(s) + w_{riv} \cdot D_{riv}(s)
\label{eq:utilidad}
\end{equation}

Donde:
\begin{align}
D_{ent}(s) &= \begin{cases}
+1 & \text{si MAX entregó (ganó)} \\
-1 & \text{si MIN entregó (ganó)} \\
0 & \text{si no hay victoria}
\end{cases} \\
C_{car}(s) &= \begin{cases}
+1 & \text{si MAX porta el recurso} \\
-1 & \text{si MIN porta el recurso} \\
0 & \text{si nadie porta}
\end{cases} \\
D_{goal}(s) &= \text{distancia Manhattan al objetivo actual} \\
D_{riv}(s) &= \begin{cases}
\min_{r \in Rivales} dist(r, pos_{portador}) & \text{si MAX porta} \\
0 & \text{si MIN porta o nadie porta}
\end{cases}
\end{align}

\section{Justificación de Cada Término}
\label{sec:justificacion-terminos}

\begin{table}[H]
\centering
\caption{Justificación de términos de la utilidad}
\begin{tabular}{llp{6cm}}
\toprule
\textbf{Término} & \textbf{Peso} & \textbf{Por qué} \\
\midrule
Entrega ($w_{ent}=10$) & 10 & Victoria/derrota es el objetivo final; debe dominar cualquier combinación de distancias \\
Carga ($w_{car}=3$) & 3 & Portar el recurso es el paso crítico hacia la victoria; vale más que distancia \\
Distancia objetivo ($w_{dist}=1$) & 1 & Cada paso de Manhattan = -1; progreso visible y monótono \\
Distancia rival ($w_{riv}=0.5$) & 0.5 & Rival cerca = riesgo de interceptación; rompe empates entre rutas equidistantes \\
\bottomrule
\end{tabular}
\end{table}
\label{tab:justificacion-utilidad}


\section{Objetivo Dinámico (Distancia al Objetivo)}
\label{sec:objetivo-dinamico}

El objetivo hacia el que se mide la distancia \textbf{cambia según quién porta el recurso}:

\begin{table}[H]
\centering
\caption{Objetivo de la distancia según portador}
\begin{tabular}{lll}
\toprule
\textbf{Situación} & \textbf{Objetivo} & \textbf{Desde dónde se mide} \\
\midrule
MAX porta & Base propia de MAX & Desde el portador \\
MIN porta & Posición del portador (perseguir) & Desde la unidad MAX más cercana \\
Nadie porta & Celda del recurso & Desde la unidad MAX más cercana \\
\bottomrule
\end{tabular}
\end{table}
\label{tab:objetivo-dinamico}


\textbf{Por qué:} Si MIN porta, MAX debe \textit{perseguir al portador} para interceptar o
forzar errores; medir distancia al recurso sería contraproducente (el recurso se mueve con
MIN). Si nadie porta, ambos bandos van al recurso.

\section{Termino Rival: Solo Cuando MAX Porta}
\label{sec:termino-rival}

\begin{equation}
D_{riv}(s) = \begin{cases}
\min_{r \in Unidades_{MIN}} dist(r, pos_{portador_{MAX}}) & \text{si MAX porta} \\
0 & \text{en otro caso}
\end{cases}
\end{equation}

\textbf{Por qué solo cuando MAX porta:} Un rival cerca del portador de MAX = amenaza de
interceptación (pierde el recurso). Si MIN porta, un rival de MAX cerca del portador de MIN
\textit{no} es una amenaza para MAX (es una oportunidad de interceptación, que ya captura la
persecución via $D_{goal}$).

\section{Pesos Configurables en JSON (RNF-03)}
\label{sec:pesos-json}

Los pesos se leen de la sección opcional \code{utilidad} del JSON:

\begin{minted}{json}
{
  "utilidad": {
    "entrega": 10,
    "recurso_cargado": 3,
    "distancia_objetivo": 1,
    "distancia_rival": 0.5
  }
}
\end{minted}

Si falta la clave, se usa el default (10/3/1/0.5). \textbf{Nada está hardcodeado} (RNF-03).

\section{Limite Conocida: Recurso en Celda Propia vs. Llevarlo Lejos}
\label{sec:limitacion-utilidad}

Con pesos 10/3/1/0.5, surge una anomalía documentada en \S\ref{sec:limitaciones-conocidas}:

\begin{itemize}
\item Recurso en celda propia de MAX (nadie porta, distancia 0): $U = -1 \cdot 0 = 0$
\item MAX porta y está a 4 pasos de su base: $U = +3 - 1\cdot4 = -1$
\item \textbf{Resultado:} Recurso en celda propia ($U=0$) se valora \textbf{mejor} que llevarlo
  a más de 3 pasos de la base ($U < 0$).
\item Además, profundidades pares aplanan valores a 0 (efecto horizonte, caso 10).
\end{itemize}

\textbf{Decisión (2026-10-17):} No cambiar defaults (invalidaría casos 7--10 validados).
Se documenta como limitación; quien necesite otro comportamiento ajusta
\code{recurso\_cargado $\ge$ 8} en el JSON de su escenario.

\section{Explicación para Sustentación: Ejemplo Caso 7}
\label{sec:ejemplo-caso7}

Escenario \code{caso\_07\_decision\_adversarial.json}: mapa 5$\times$7, coste 1, recurso en
(0,1), A1 en (1,0), A2 en (0,2), B1 en (1,2), B2 en (2,4), A mueve.

\begin{enumerate}
\item \textbf{Profundidad 1 (A mueve, evalúa hoja):}
  \begin{itemize}
  \item A2$\to$(0,1) agarra recurso: $U = +3 - 1\cdot1 = 2$ (distancia a base A = 1) + rival term
  \item Valor $\approx$ 3 $\to$ elige agarrar
  \end{itemize}
\item \textbf{Profundidad 2 (B responde):}
  \begin{itemize}
  \item B1 se mueve adyacente a A2 (intercepta): recurso cae, $U$ cae a 0
  \item Valor 0 $\to$ sigue agarrando (empate técnico, tie-break orden)
  \end{itemize}
\item \textbf{Profundidad 3 (A vuelve a mover):}
  \begin{itemize}
  \item A1$\to$(1,1) \textit{protege} (bloquea a B1 antes de que A2 agarre)
  \item B1 no puede interceptar $\to$ A2 agarra y mantiene $\to$ $U = 3$
  \end{itemize}
\item \textbf{Profundidad 4:} A1$\to$(0,0) (entra a base sin recurso) valor 0; B responde...
\end{enumerate}

\textbf{Conclusión:} La utilidad hace que el sistema \textbf{aprenda a proteger} antes de
agarrar, porque la interceptación del rival hace que el valor de agarrar directo caiga a 0.
% 10_experimentacion.tex - Capítulo 10: Experimentación y Resultados
\chapter{Metodología Experimental y Resultados}
\label{ch:experimentos}

\section{Metodología}
\label{sec:metodologia}

\screenshotfigure{cli-resultado}{0.90}{Ejecución de una búsqueda individual con métricas y verificaciones RF-42.}

\screenshotfigure{cli-sweep}{0.90}{Sweep de algoritmos con exportación de resultados en CSV y JSON.}

\begin{itemize}
\item \textbf{Escenarios:} Apéndice A oficial (\code{apendice\_a.json}) + casos del equipo en
  \code{scenarios/} (\code{caso\_01}..\code{caso\_07}, \code{prueba\_15x15}, \code{prueba\_20x20},
  \code{escenario\_03x03}..\code{escenario\_30x30}).
\item \textbf{Valores esperados:} Derivados \textbf{a mano del mapa} \textbf{antes} de ejecutar
  el código; aserciones en \code{tests/test\_cases\_f*.py}.
\item \textbf{Determinismo:} Métricas de conteo (generados, expandidos, frontera, costo, longitud)
  son idénticas en cada ejecución. Solo \code{time\_seconds} varía (tiempo de ejecución).
\item \textbf{Reproducibilidad:} Cada número se regenera con los comandos del capítulo \ref{ch:reproducibilidad}
  (\texttt{user-manual.md} §6). Archivos brutos en \texttt{docs/results/}.
\end{itemize}

\section{Casos 1 y 2: No Informada y Costos (F2)}
\label{sec:resultados-f2}

\subsection{Caso 1: BFS vs DFS (Mismo costo, distinta exploración)}
\label{sec:caso1}

\textbf{Setup:} Grid 20$\times$20, todo \code{camino} (costo 1). A1 en (2,2), meta (2,6).

\begin{table}[H]
\centering
\caption{Caso 1: BFS vs DFS}
\begin{tabular}{lccccccc}
\toprule
Algoritmo & Éxito & Costo & Movs & Generados & Expandidos & Max Frontera & Match \\
\midrule
BFS & Sí & 4 & 4 & 39 & 29 & 11 & Sí \\
DFS & Sí & 42 & 42 & 152 & 94 & 58 & Sí \\
\bottomrule
\end{tabular}
\end{table}
\label{tab:caso1}

\textbf{Conclusión:} Ambos encuentran solución (completitud en grafo finito), pero DFS explora
un camino 10$\times$ más largo porque expande ``abajo'' primero. BFS garantiza mínimos
movimientos.

\subsection{Caso 2: Ruta Corta vs. Ruta Económica}
\label{sec:experimento-caso2-detalle}

\textbf{Setup:} Banda de pantano (costo 7) en columnas 8--10. A1 (5,5), meta (5,12).

\begin{table}[H]
\centering
\caption{Caso 2: BFS vs UCS}
\begin{tabular}{lccccccc}
\toprule
Algoritmo & Éxito & Costo & Movs & Generados & Expandidos & Max Frontera & Match \\
\midrule
BFS & Sí & 25 & 7 & 126 & 104 & 23 & Sí \\
UCS & Sí & 17 & 17 & 269 & 239 & 33 & Sí \\
\bottomrule
\end{tabular}
\end{table}
\label{tab:caso2-detalle}

\textbf{Conclusión:} La ruta con menos movimientos (BFS, 7 movs) cuesta 25; UCS rodea el
pantano (17 movs) y paga 17. \textbf{La ruta con menos movimientos no es la más barata.}

\section{Casos 3, 4, 5: Informada y Memoria Limitada (F3)}
\label{sec:resultados-f3}

\subsection{Caso 3: Obstáculo y Heurística}
\label{sec:caso3}

Muro vertical en columna 8 (filas 0--15). A1 (5,5), meta (5,12) $\to$ debe rodear por abajo
(fila 16). Distancia real = 29 movimientos.

\begin{table}[H]
\centering
\caption{Caso 3: Expansiones -- h2 ve el muro, h1 no}
\begin{tabular}{lccccc}
\toprule
Algoritmo & Costo & Movs & Generados & Expandidos & Max Frontera \\
\midrule
UCS & 29 & 29 & 307 & 295 & 17 \\
A* h1 & 29 & 29 & 194 & 166 & 28 \\
A* h2 & 29 & 29 & 122 & 84 & 38 \\
\bottomrule
\end{tabular}
\end{table}
\label{tab:caso3-expansiones}

h2 expande \textbf{2$\times$ menos} que h1 porque ``ve'' el muro y no explora el lado
equivocado.

\subsection{Caso 4: Comparación de Heurísticas}
\label{sec:caso4}

Muro en columna 9 (filas 3--16). Hueco arriba (filas 0--2) = pantano (costo 3); hueco abajo
= camino (costo 1). Óptimo = 22 movimientos, costo 24 (por abajo).

\begin{table}[H]
\centering
\caption{Caso 4: A* con h1 vs h2}
\begin{tabular}{lccccc}
\toprule
Algoritmo & Costo & Movs & Generados & Expandidos & Max Frontera \\
\midrule
UCS & 24 & 22 & 346 & 332 & 20 \\
A* h1 & 24 & 22 & 230 & 186 & 44 \\
A* h2 & 24 & 22 & 164 & 115 & 49 \\
\bottomrule
\end{tabular}
\end{table}
\label{tab:caso4}

h2 expande menos que h1 (ve el muro), pero \textbf{no ve el pantano}: el atajo por arriba
es más corto en movimientos pero cruza pantano (costo 3). h2 lo explora igual que h1.

\subsection{Caso 5: Beam Search y k}
\label{sec:caso5}

Mapa 7$\times$9, hierba (2) y bosque (4). Óptimo: costo 28, 14 movimientos (UCS).

\begin{table}[H]
\centering
\caption{Caso 5: Beam Search con k variable}
\begin{tabular}{ccccccc}
\toprule
k & Éxito & Costo & Movs & Generados & Expandidos & Max Frontera \\
\midrule
1 & No & -- & -- & 9 & 5 & 3 \\
2 & Sí & 30 & 16 & 43 & 29 & 6 \\
4 & Sí & 28 & 14 & 50 & 41 & 6 \\
8 & Sí & 28 & 14 & 46 & 44 & 6 \\
\bottomrule
\end{tabular}
\end{table}
\label{tab:caso5-beam}

\begin{itemize}
\item $k=1$: incompleto, poda la única rama prometedora $\to$ falla.
\item $k=2$: sobrevive pero mantiene ruta subóptima (30 > 28).
\item $k=4,8$: recuperan óptimo. $k=8$ expande más que $k=4$ para mismo camino (punto de
  rendimientos decrecientes).
\end{itemize}

\section{Casos 6--10: Adversarial (F4)}
\label{sec:resultados-f4}

Ver Cap. \ref{ch:adversarial}, \S\ref{sec:casos7-10}. Resumen:

\begin{table}[H]
\centering
\caption{Resumen casos 6--10}
\begin{tabular}{lll}
\toprule
Caso & Descripción & Resultado \\
\midrule
6 & Estado compuesto & 3 estados en misma posición (1,0) $\to$ 3 objetos distintos \checkmark \\
7 & Decisión adversarial & Prof 1/2 agarran; prof 3 protege (A1$\to$(1,1)) \checkmark \\
8 & Minimax vs $\alpha$-$\beta$ & Mismo valor/acción; $\alpha$-$\beta$ 11$\times$ menos nodos (d=4) \checkmark \\
9 & Orden de acciones & Mismo valor, distinta poda (d=3: 142 vs 191 gen) \checkmark \\
10 & Profundidad & Impares=guarda (3), pares=aplana a 0 (efecto horizonte) \checkmark \\
\bottomrule
\end{tabular}
\end{table}
\label{tab:casos6-10-resumen}


\section{Rendimiento (G13)}
\label{sec:rendimiento-detalle}

Escenario \code{prueba\_15x15.json} (15$\times$15, pantanos, muro, 3 unidades/bando).
Profundidad = acciones de un bando.

\begin{table}[H]
\centering
\caption{Rendimiento Minimax vs Alfa-Beta (15$\times$15, 3 unidades/bando)}
\begin{tabular}{lcc}
\toprule
Algoritmo & Profundidad & Tiempo \\
\midrule
Alpha-beta & 5 & 0.04 s \\
Alpha-beta & 6 & 0.18 s \\
Alpha-beta & 7 & 0.86 s \\
Minimax & 5 & 1.2 s \\
Minimax & 6 & 15.7 s (supera meta 5 s) \\
\bottomrule
\end{tabular}
\end{table}
\label{tab:rendimiento-15x15}


Con 5 unidades/bando: $\alpha$-$\beta$ d=6 $\to$ 0.85 s. En 20$\times$20 (3 unidades):
$\alpha$-$\beta$ d=2..6 $\le$ 0.22 s.

\textbf{Conclusión:} Dentro de la meta de 5 s, $\alpha$-$\beta$ hasta profundidad 7; Minimax
simple hasta profundidad 5. Medido en laptop del equipo (Windows, Python 3.14).
% 10b_resultados.tex - Capítulo 10b: Tablas completas de resultados (Casos 1-10)
\chapter{Resultados Completos: Tablas Esperado vs. Obtenido}
\label{ch:resultados}

Este capítulo presenta las tablas completas de los 10 casos mínimos de demostración con
valores esperados (derivados a mano) vs. obtenidos (ejecución real con \code{--verify}).
Todos los \textbf{Match = yes}.

\section{Caso 1: BFS vs DFS}
\label{sec:tabla-caso1}

\begin{table}[H]
\centering
\caption{Caso 1 -- BFS vs DFS (mapa 20$\times$20, todo coste 1, A1 (2,2) $\to$ (2,6))}
\begin{tabular}{lccccccc}
\toprule
Algoritmo & Éxito & Costo & Movs & Generados & Expandidos & Max Frontera & Match \\
\midrule
BFS & Sí & 4 & 4 & 39 & 29 & 11 & Sí \\
DFS & Sí & 42 & 42 & 152 & 94 & 58 & Sí \\
\bottomrule
\end{tabular}
\end{table}
\label{tab:caso1-full}


\section{Caso 2: BFS vs UCS}
\label{sec:tabla-caso2}

\begin{table}[H]
\centering
\caption{Caso 2 -- BFS vs UCS (banda pantano coste 7, A1 (5,5) $\to$ (5,12))}
\begin{tabular}{lccccccc}
\toprule
Algoritmo & Éxito & Costo & Movs & Generados & Expandidos & Max Frontera & Match \\
\midrule
BFS & Sí & 25 & 7 & 126 & 104 & 23 & Sí \\
UCS & Sí & 17 & 17 & 269 & 239 & 33 & Sí \\
\bottomrule
\end{tabular}
\end{table}
\label{tab:caso2-full}


\section{Caso 3: Obstáculo y Heurística}
\label{sec:tabla-caso3}

\begin{table}[H]
\centering
\caption{Caso 3 -- Muro vertical, A* h1 vs h2 (A1 (5,5) $\to$ (5,12), muro col 8)}
\begin{tabular}{lcccccc}
\toprule
Algoritmo & Costo & Movs & Generados & Expandidos & Max Frontera & Match \\
\midrule
UCS & 29 & 29 & 307 & 295 & 17 & Sí \\
A* h1 & 29 & 29 & 194 & 166 & 28 & Sí \\
A* h2 & 29 & 29 & 122 & 84 & 38 & Sí \\
\bottomrule
\end{tabular}
\end{table}
\label{tab:caso3-full}


\section{Caso 4: Comparación de Heurísticas}
\label{sec:tabla-caso4}

\begin{table}[H]
\centering
\caption{Caso 4 -- Muro con hueco pantano vs. libre (A1 (9,5) $\to$ (9,13))}
\begin{tabular}{lcccccc}
\toprule
Algoritmo & Costo & Movs & Generados & Expandidos & Max Frontera & Match \\
\midrule
UCS & 24 & 22 & 346 & 332 & 20 & Sí \\
A* h1 & 24 & 22 & 230 & 186 & 44 & Sí \\
A* h2 & 24 & 22 & 164 & 115 & 49 & Sí \\
\bottomrule
\end{tabular}
\end{table}
\label{tab:caso4-full}


\section{Caso 5: Beam Search k=1,2,4,8}
\label{sec:tabla-caso5}

\begin{table}[H]
\centering
\caption{Caso 5 -- Beam Search k=1,2,4,8 (7$\times$9, óptimo costo 28, 14 movs)}
\begin{tabular}{cccccccc}
\toprule
k & Éxito & Costo & Movs & Generados & Expandidos & Max Frontera & Match \\
\midrule
1 & No & -- & -- & 9 & 5 & 3 & Sí \\
2 & Sí & 30 & 16 & 43 & 29 & 6 & Sí \\
4 & Sí & 28 & 14 & 50 & 41 & 6 & Sí \\
8 & Sí & 28 & 14 & 46 & 44 & 6 & Sí \\
\bottomrule
\end{tabular}
\end{table}
\label{tab:caso5-full}


\section{Caso 6: Estado Compuesto}
\label{sec:tabla-caso6}

\begin{table}[H]
\centering
\caption{Caso 6 -- Misma posición (1,0), estados distintos}
\begin{tabular}{lllll}
\toprule
Estado & Pos A1 & Carrier & Turn & Efecto Move A1 $\to$ (0,0) \\
\midrule
Libre & (1,0) & nadie & A & entra base, no gana \\
Portando & (1,0) & A1 & A & \textbf{gana A} \\
Turno rival & (1,0) & nadie & B & sin acciones legales \\
\bottomrule
\end{tabular}
\end{table}
\label{tab:caso6-full}


Verificación: \code{len(set(states)) == 3} (3 objetos distintos), \code{winner} correcto,
acciones legales coherentes. Test: \code{tests/test\_case\_6.py}.

\section{Caso 7: Decisión Adversarial}
\label{sec:tabla-caso7}

\begin{table}[H]
\centering
\caption{Caso 7 -- Minimax profundidades 1--3 (mapa 5$\times$7, recurso (0,1))}
\begin{tabular}{lccc}
\toprule
Profundidad & Acción esperada & Valor esperado & Match \\
\midrule
1 & A2$\to$(0,1) & 3 & Sí \\
2 & A2$\to$(0,1) & 0 & Sí \\
3 & A1$\to$(1,1) & 3 & Sí \\
\bottomrule
\end{tabular}
\end{table}
\label{tab:caso7-full}


\section{Caso 8: Minimax vs Alfa-Beta}
\label{sec:tabla-caso8}

\begin{table}[H]
\centering
\caption{Caso 8 -- Profundidad 4, orden default}
\begin{tabular}{lcccccc}
\toprule
Algoritmo & Acción & Valor & Generados & Evaluados & Podados & Match \\
\midrule
Minimax & A1$\to$(0,0) & 0 & 4,826 & 4,296 & 0 & Sí \\
Alpha-beta & A1$\to$(0,0) & 0 & 431 & 299 & 757 & Sí \\
\bottomrule
\end{tabular}
\end{table}
\label{tab:caso8-full}


\section{Caso 9: Orden de Acciones}
\label{sec:tabla-caso9}

\begin{table}[H]
\centering
\caption{Caso 9 -- Profundidad 3, orden default vs reversed}
\begin{tabular}{lcccccc}
\toprule
Orden & Acción & Valor & Generados & Evaluados & Podados & Match \\
\midrule
default & A1$\to$(1,1) & 3 & 142 & 110 & 106 & Sí \\
reversed & A1$\to$(1,1) & 3 & 191 & 148 & 145 & Sí \\
\bottomrule
\end{tabular}
\end{table}
\label{tab:caso9-full}


\section{Caso 10: Profundidad}
\label{sec:tabla-caso10}

\begin{table}[H]
\centering
\caption{Caso 10 -- Alfa-beta default order, profundidades 1--5}
\begin{tabular}{lcccc}
\toprule
Profundidad & Acción & Valor & Evaluados & Match \\
\midrule
1 & A2$\to$(0,1) & 3 & 7 & Sí \\
2 & A2$\to$(0,1) & 0 & 24 & Sí \\
3 & A1$\to$(1,1) & 3 & 110 & Sí \\
4 & A1$\to$(0,0) & 0 & 299 & Sí \\
5 & A1$\to$(1,1) & 3 & 1,583 & Sí \\
\bottomrule
\end{tabular}
\end{table}
\label{tab:caso10-full}


(Con Minimax: evaluados = 7 / 64 / 467 / 4,296 / 30,994)
% 11_verificacion_testing.tex - Capítulo 11: Verificación y Testing
\chapter{Verificación, Testing y Calidad}
\label{ch:verificacion}

\section{Verificación de Soluciones (RF-42)}
\label{sec:verificacion-rf42}

El enunciado exige 6 chequeos (RF-42). Todos implementados en \code{experiments/verifier.py}:

\begin{table}[H]
\centering
\caption{6 Checks de verificación (RF-42)}
\begin{tabular}{llp{5.5cm}}
\toprule
\# & Check & Descripción \\
\midrule
1 & \textit{valid moves} & Cada paso del camino es acción legal según \code{legal\_actions} \\
2 & \textit{no obstacles} & Ninguna celda del camino es no transitable \\
3 & \textit{cost matches} & Costo reportado = suma de costos al re-ejecutar (\code{math.isclose}) \\
4 & \textit{goal reached} & Última celda del camino = \code{problem.goal} \\
5 & \textit{no invalid states} & Estados rejugados: 1 unidad/celda, todas transitables \\
6 & \textit{Minimax legal} & Acciones consideradas por Minimax/$\alpha$-$\beta$ son legales; acción elegida es legal \\
\bottomrule
\end{tabular}
\end{table}
\label{tab:checks-rf42}


Implementación:
\begin{itemize}
\item Checks 1--5: \code{verify\_search(problem, result)} re-ejecuta el camino paso a paso con
  \code{problem.expand} y valida cada transición.
\item Check 6: \code{verify\_adversarial(board, state, result)} usa \code{record=True} en
  Minimax/$\alpha$-$\beta$ para obtener \code{considered} y valida legalidad de cada acción
  considerada y de la elegida.
\end{itemize}

Todos los reportes en \texttt{docs/results/} se generaron con \code{--verify} y \textbf{todos
los checks pasan}.

\section{Testing Automatizado}
\label{sec:testing}

\screenshotfigure{test-resultados}{0.90}{Ejecución de pytest, ruff y mypy con resultados satisfactorios.}

\begin{table}[H]
\centering
\caption{Suite de tests (281 tests, 0 skips)}
\begin{tabular}{ll}
\toprule
\textbf{Archivo} & \textbf{Qué prueba} \\
\midrule
\code{test\_domain\_state.py} & \code{GameState}, \code{Position}, \code{Board}, inmutabilidad \\
\code{test\_domain\_rules.py} & \code{legal\_actions}, \code{apply}, \code{successors}, \code{winner}, \code{is\_terminal} \\
\code{test\_scenario\_validation.py} & 9 validaciones RF-35 + extra (célula compartida), carga atómica \\
\code{test\_bfs.py}, \code{test\_dfs.py}, \code{test\_ucs.py} & Métricas, optimalidad, convenciones \\
\code{test\_search\_common.py} & \code{make\_problem}, \code{reconstruct}, \code{\_normalize} \\
\code{test\_heuristics.py} & Escalado \code{min\_step\_cost}, dominancia h2$\ge$h1, celdas inalcanzables \\
\code{test\_astar.py} & Optimalidad, equivalencia con UCS, reapertura \\
\code{test\_beam.py} & Límite k, fallo k=1, determinismo \\
\code{test\_adversarial.py} & Minimax, $\alpha$-$\beta$, métricas RF-30, mismo decisión (RF-31) \\
\code{test\_case\_6.py} & Caso 6: estado compuesto, 3 estados misma posición \\
\code{test\_cases\_f2.py}, \code{test\_cases\_f3.py}, \code{test\_cases\_f4.py} & Casos 1--10 esperado vs obtenido \\
\code{test\_utility.py} & Pesos, objetivo dinámico, término rival, terminales \\
\code{test\_verifier.py} & 6 checks RF-42 (búsqueda + adversarial) \\
\code{test\_runner.py}, \code{test\_sweep.py} & CLI, sweep combinatorio, export CSV/JSON \\
\code{test\_gui\_app.py} & Ventana, dibujo, clics, auto-mode \\
\code{test\_runner.py} & CLI, modos, sweep, export \\
\bottomrule
\end{tabular}
\end{table}
\label{tab:tests}


\textbf{Comando:} \code{uv run pytest} (281 tests, 0 skips).
\textbf{Calidad:} \code{uv run ruff check .} + \code{uv run mypy src --strict} $\to$ \textbf{verde}.

\section{Cobertura de Requisitos por Test}
\label{sec:cobertura-rf}

\begin{table}[H]
\centering
\caption{Cobertura de RF por archivo de test}
\begin{tabular}{ll}
\toprule
\textbf{RF} & \textbf{Test que lo cubre} \\
\midrule
RF-03, RF-04, RF-37 & \code{test\_scenario\_validation.py} (costos desde JSON) \\
RF-05, RF-06, RF-07, RF-08, RF-09 & \code{test\_domain\_rules.py} \\
RF-10, RF-11, RF-12, RF-13 & \code{test\_domain\_state.py}, \code{test\_case\_6.py} \\
RF-14, RF-15, RF-16, RF-19 & \code{test\_bfs.py}, \code{test\_dfs.py}, \code{test\_ucs.py} \\
RF-20, RF-21, RF-22, RF-23, RF-24 & \code{test\_astar.py}, \code{test\_beam.py}, \code{test\_heuristics.py} \\
RF-25--33 & \code{test\_adversarial.py}, \code{test\_cases\_f4.py} \\
RF-34, RF-35 & \code{test\_scenario\_validation.py} \\
RF-42 (1-5) & \code{test\_verifier.py} (búsqueda) \\
RF-42 (6) & \code{test\_verifier.py} (adversarial), \code{test\_adversarial.py} \\
\bottomrule
\end{tabular}
\end{table}
\label{tab:cobertura-rf}


\section{Calidad de Código}
\label{sec:calidad}

\begin{itemize}
\item \code{uv run ruff check .} $\to$ \textbf{0 errores, 0 warnings}
\item \code{uv run ruff format --check .} $\to$ \textbf{formato correcto}
\item \code{uv run mypy src --strict} $\to$ \textbf{0 errores} (tipado estricto)
\item Docstrings con: Purpose, Preconditions, Postconditions, Complexity, AI usage, AI intervention, Student validation (sección 9.2 enunciado)
\item \textbf{Validación del equipo:} Revisión cruzada realizada por los integrantes
\end{itemize}

\section{Integración Continua (Manual)}
\label{sec:ci}

No hay CI/CD automatizado configurado, pero el flujo de trabajo manual antes de cada commit:
\begin{enumerate}
\item \code{uv run pytest}
\item \code{uv run ruff check .}
\item \code{uv run ruff format --check .}
\item \code{uv run mypy src}
\item Si todo verde $\to$ commit
\end{enumerate}
% 11b_casos_demostracion.tex - Capítulo 11b: Casos Mínimos de Demostración (detalle)
\chapter{Casos Mínimos de Demostración: Detalle Completo}
\label{ch:casos-demo}

Este capítulo presenta los 10 casos mínimos del enunciado (Sección 13) con sus tablas
completas, mapas y comandos de reproducción. Cada caso incluye: \textbf{Setup}, \textbf{Expectativa},
\textbf{Tabla esperado vs. obtenido}, \textbf{Comando de reproducción}.

Todos los valores esperados fueron \textbf{derivados a mano del mapa antes de ejecutar el código}
y son asertados por los tests automatizados.

\screenshotfigure{gui-principal}{0.95}{Vista principal de la interfaz gráfica con el escenario cargado.}

\screenshotfigure{gui-controles}{0.80}{Controles de algoritmo y parámetros de ejecución.}

\screenshotfigure{gui-resultados}{0.85}{Resultados de una ejecución mostrados en la interfaz gráfica.}

\screenshotfigure{gui-jugadas}{0.80}{Traza e historial de acciones de una partida.}

\screenshotfigure{gui-automatico}{0.85}{Interfaz durante una partida automática IA contra IA.}

\section{Caso 1: BFS frente a DFS}
\label{sec:caso1-detalle}

\textbf{Archivo:} \code{scenarios/caso\_01\_bfs\_vs\_dfs.json}
\textbf{Mapa:} 20$\times$20, todo \code{camino} (costo 1).
\textbf{Unidades:} A1 en (2,2), meta (2,6).

\textbf{Expectativa:} Ambos encuentran solución pero exploran el espacio de manera distinta.
BFS retorna el camino recto (4 movs, coste 4). DFS expande ``abajo'' primero y llega por un
camino mucho más largo (42 movs, coste 42).

\begin{table}[H]
\centering
\caption{Caso 1 -- BFS vs DFS}
\begin{tabular}{lccccccc}
\toprule
Algoritmo & Éxito & Costo & Movs & Generados & Expandidos & Max Frontera & Match \\
\midrule
BFS & Sí & 4 & 4 & 39 & 29 & 11 & Sí \\
DFS & Sí & 42 & 42 & 152 & 94 & 58 & Sí \\
\bottomrule
\end{tabular}
\end{table}

\textbf{Comando:} \code{uv run python -m tactical\_grid scenarios/caso\_01\_bfs\_vs\_dfs.json --algorithm BFS,DFS --verify}

\section{Caso 2: Ruta Corta vs. Económica}
\label{sec:caso2-detalle}

\textbf{Archivo:} \code{scenarios/caso\_02\_bfs\_vs\_ucs.json}
\textbf{Mapa:} 20$\times$20, banda de pantano (coste 7) en columnas 8--10, filas 0--9.
A1 en (5,5), meta (5,12).

\textbf{Expectativa:} BFS cruza la banda: 7 movs, coste $4\times1 + 3\times7 = 25$.
UCS rodea por fila 10: 17 movs, coste 17.

\begin{table}[H]
\centering
\caption{Caso 2 -- BFS vs UCS}
\begin{tabular}{lccccccc}
\toprule
Algoritmo & Éxito & Costo & Movs & Generados & Expandidos & Max Frontera & Match \\
\midrule
BFS & Sí & 25 & 7 & 126 & 104 & 23 & Sí \\
UCS & Sí & 17 & 17 & 269 & 239 & 33 & Sí \\
\bottomrule
\end{tabular}
\end{table}

\textbf{Conclusión:} La ruta con menos movimientos \textbf{no} es la más barata.

\section{Caso 3: Obstáculo y Heurística}
\label{sec:caso3-detalle}

\textbf{Archivo:} \code{scenarios/caso\_03\_obstaculo\_heuristica.json}
\textbf{Mapa:} Muro vertical en columna 8 (filas 0--15). A1 (5,5), meta (5,12).
Distancia real = 29 (rodear por fila 16).

\textbf{Expectativa:} Manhattan (h1) es engañado por el muro; obstacle-aware (h2) no.
Expansiones: h2 $<$ h1 $<$ UCS.

\begin{table}[H]
\centering
\caption{Caso 3 -- UCS, A* h1, A* h2}
\begin{tabular}{lcccccc}
\toprule
Algoritmo & Costo & Movs & Generados & Expandidos & Max Frontera & Match \\
\midrule
UCS & 29 & 29 & 307 & 295 & 17 & Sí \\
A* h1 & 29 & 29 & 194 & 166 & 28 & Sí \\
A* h2 & 29 & 29 & 122 & 84 & 38 & Sí \\
\bottomrule
\end{tabular}
\end{table}

\textbf{Conclusión:} h2 ve el muro y no explora el lado equivocado; h1 explora ambos lados.

\section{Caso 4: Comparación de Heurísticas}
\label{sec:caso4-detalle}

\textbf{Archivo:} \code{scenarios/caso\_04\_comparacion\_heuristicas.json}
\textbf{Mapa:} Muro en columna 9 (filas 3--16). Hueco arriba (filas 0--2) = pantano (coste 3).
Hueco abajo = camino (coste 1). Óptimo = 22 movs, coste 24 (por abajo).

\textbf{Expectativa:} h2 expande menos que h1 (ve el muro), pero no ve el pantano.

\begin{table}[H]
\centering
\caption{Caso 4 -- UCS, A* h1, A* h2}
\begin{tabular}{lcccccc}
\toprule
Algoritmo & Costo & Movs & Generados & Expandidos & Max Frontera & Match \\
\midrule
UCS & 24 & 22 & 346 & 332 & 20 & Sí \\
A* h1 & 24 & 22 & 230 & 186 & 44 & Sí \\
A* h2 & 24 & 22 & 164 & 115 & 49 & Sí \\
\bottomrule
\end{tabular}
\end{table}

\textbf{Conclusión:} h2 ayuda (expande menos que h1) pero no ve el pantano $\to$ sigue
explorando el atajo costoso.

\section{Caso 5: Beam Search}
\label{sec:caso5-detalle}

\textbf{Archivo:} \code{scenarios/caso\_05\_beam\_k.json}
\textbf{Mapa:} 7$\times$9, hierba (coste 2) y bosque (coste 4). Óptimo: coste 28, 14 movs (UCS).

\textbf{Expectativa:} $k=1$ poda la única rama buena $\to$ falla. $k=2$ sobrevive pero mantiene
ruta subóptima (30, 16 movs). $k=4,8$ recuperan óptimo.

\begin{table}[H]
\centering
\caption{Caso 5 -- Beam Search k=1,2,4,8}
\begin{tabular}{cccccccc}
\toprule
k & Éxito & Costo & Movs & Generados & Expandidos & Max Frontera & Match \\
\midrule
1 & No & -- & -- & 9 & 5 & 3 & Sí \\
2 & Sí & 30 & 16 & 43 & 29 & 6 & Sí \\
4 & Sí & 28 & 14 & 50 & 41 & 6 & Sí \\
8 & Sí & 28 & 14 & 46 & 44 & 6 & Sí \\
\bottomrule
\end{tabular}
\end{table}

\textbf{Conclusión:} Restringir alternativas \textbf{puede descartar un camino que después
habría sido conveniente} (RF-24). $k=1$ falla; $k=2$ sobrevive pero subóptimo; $k\ge4$
recupera óptimo.

\section{Caso 6: Estado Compuesto}
\label{sec:caso6-detalle}

\textbf{Archivo:} \code{scenarios/caso\_07\_decision\_adversarial.json} (mapa 5$\times$7).
A1 en (1,0), base A en (0,0).

\textbf{Expectativa:} La misma posición física (1,0) representa estados distintos según
información adicional del juego:

\begin{table}[H]
\centering
\caption{Caso 6 -- Estado compuesto}
\begin{tabular}{lllll}
\toprule
Estado & Pos A1 & Carrier & Turn & Efecto Move A1 $\to$ (0,0) \\
\midrule
Libre (inicial) & (1,0) & nadie & A & entra base, no gana \\
Portando & (1,0) & A1 & A & \textbf{gana A} \\
Turno rival & (1,0) & nadie & B & A1 sin acciones legales \\
\bottomrule
\end{tabular}
\end{table}

Test \code{tests/test\_case\_6.py}: \code{len(set(states)) == 3}, \code{winner} correcto,
acciones legales coherentes. Un dict indexado por estado mantiene 3 entradas; por posición
solo 1.

\section{Caso 7: Decisión Adversarial}
\label{sec:caso7-detalle}

\textbf{Archivo:} \code{scenarios/caso\_07\_decision\_adversarial.json}
Mapa 5$\times$7, coste 1, recurso (0,1), A1 (1,0), A2 (0,2), B1 (1,2), B2 (2,4), A mueve.
Pesos default: 10/3/1/0.5.

\textbf{Expectativa:} La mejor decisión no es ir directo al recurso.

\begin{table}[H]
\centering
\caption{Caso 7 -- Minimax profundidades 1--3}
\begin{tabular}{lccc}
\toprule
Profundidad & Acción esperada & Valor esperado & Match \\
\midrule
1 & A2$\to$(0,1) & 3 & Sí \\
2 & A2$\to$(0,1) & 0 & Sí \\
3 & A1$\to$(1,1) & 3 & Sí \\
\bottomrule
\end{tabular}
\end{table}

\textbf{Análisis:} Prof 1/2 agarran directo (valor 3/0). Prof 3 descubre que mover A1$\to$(1,1)
primero bloquea a B1 y permite que A2 agarre seguro: $U = 3 - 1 + 0.5\times2 = 3$.

\section{Caso 8: Minimax vs Alfa-Beta}
\label{sec:caso8-detalle}

\begin{table}[H]
\centering
\caption{Caso 8 -- Profundidad 4, orden default}
\begin{tabular}{lcccccc}
\toprule
Algoritmo & Acción & Valor & Generados & Evaluados & Podados & Match \\
\midrule
Minimax & A1$\to$(0,0) & 0 & 4,826 & 4,296 & 0 & Sí \\
Alpha-beta & A1$\to$(0,0) & 0 & 431 & 299 & 757 & Sí \\
\bottomrule
\end{tabular}
\end{table}

\textbf{Conclusión:} Misma decisión, mismo valor; $\alpha$-$\beta$ genera \textbf{11$\times$ menos
nodos} y poda 757 ramas.

\section{Caso 9: Orden de Acciones}
\label{sec:caso9-detalle}

\begin{table}[H]
\centering
\caption{Caso 9 -- Profundidad 3, default vs reversed}
\begin{tabular}{lcccccc}
\toprule
Orden & Acción & Valor & Generados & Evaluados & Podados & Match \\
\midrule
default & A1$\to$(1,1) & 3 & 142 & 110 & 106 & Sí \\
reversed & A1$\to$(1,1) & 3 & 191 & 148 & 145 & Sí \\
\bottomrule
\end{tabular}
\end{table}

Misma decisión (RF-31), \textbf{poda distinta} (RF-32). En profundidad 4, orden \code{reversed}
cambia el tie-break entre valores iguales (0) y produce acción diferente con mismo valor.

\section{Caso 10: Profundidad}
\label{sec:caso10-detalle}

\begin{table}[H]
\centering
\caption{Caso 10 -- Alfa-beta default, profundidades 1--5}
\begin{tabular}{lcccc}
\toprule
Profundidad & Acción & Valor & Evaluados & Match \\
\midrule
1 & A2$\to$(0,1) & 3 & 7 & Sí \\
2 & A2$\to$(0,1) & 0 & 24 & Sí \\
3 & A1$\to$(1,1) & 3 & 110 & Sí \\
4 & A1$\to$(0,0) & 0 & 299 & Sí \\
5 & A1$\to$(1,1) & 3 & 1,583 & Sí \\
\bottomrule
\end{tabular}
\end{table}

(Minimax evaluados: 7 / 64 / 467 / 4,296 / 30,994)

\textbf{Efecto horizonte:} Profundidades impares terminan en turno de A (guarda, valor 3).
Pares terminan en turno de B (aplana a 0, tie-break por orden).
% 12_formato_json.tex - Capítulo 12: Formato JSON y Validación
\chapter{Formato JSON de Escenarios y Validación Atómica}
\label{ch:json}

\section{Contrato de Interoperabilidad}
\label{sec:contrato-json}

\screenshotfigure{json-escenario}{0.90}{Archivo JSON de escenario utilizado por la aplicación.}

\screenshotfigure{json-validacion}{0.85}{Validación de un escenario JSON.}

El formato JSON es el \textbf{contrato de interoperabilidad} (Sección 7 del enunciado):
cualquier archivo válido que respete el formato debe cargarse \textbf{sin modificar el código
fuente}. El profesor entregará al menos un escenario no visto en la sustentación.

\section{Campos Obligatorios (RF-35)}
\label{sec:campos-obligatorios}

\begin{longtable}{p{3cm}p{3cm}p{8cm}}
\toprule
\textbf{Campo} & \textbf{Tipo} & \textbf{Descripción} \\
\midrule
\endfirsthead
\toprule
\textbf{Campo} & \textbf{Tipo} & \textbf{Descripción} \\
\midrule
\endhead
\midrule
\multicolumn{3}{r}{\textit{Continúa}} \\
\endfoot
\bottomrule
\endlastfoot
\code{version} & string & Versión del formato (ej. \code{"1.0"}) \\
\code{mapa.filas} & int $>0$ & Filas del grid \\
\code{mapa.columnas} & int $>0$ & Columnas del grid \\
\code{tipos\_terreno} & object & Catálogo: \code{nombre $\to$ \{costo, transitable\}} \\
\code{terreno} & string[][] & Matriz \code{filas $\times$ columnas} con nombres de terreno \\
\code{bases.A} & \code{\{fila, columna\}} & Posición base A \\
\code{bases.B} & \code{\{fila, columna\}} & Posición base B \\
\code{recurso} & \code{\{fila, columna\}} & Posición recurso (si nadie lo transporta) \\
\code{unidades} & array & Lista: \code{id, bando, tipo, fila, columna} \\
\code{turno} & \code{"A"} | \code{"B"} & Bando que debe actuar \\
\code{juego.portador\_recurso} & string | \code{null} & id unidad que porta, o \code{null} \\
\code{juego.profundidad\_maxima\_minimax} & int (opcional) & Profundidad para Minimax \\
\code{juego.limite\_turnos} & int (opcional) & Límite turnos (default 100) \\
\code{juego.variante} & string (opcional) & \code{base} o \code{dinamica} \\
\code{juego.turnos\_recuperacion} & int (opcional) & Recuperación propia de unidades dinámicas \\
\code{juego.inmunidad\_robo} & int (opcional) & Inmunidad del portador tras un robo \\
\code{utilidad} & object (opcional) & Pesos: entrega, recurso\_cargado, distancia\_objetivo, distancia\_rival \\
\code{prueba} & object (opcional) & Configuración de prueba concreta \\
\bottomrule
\caption{Campos del formato JSON (obligatorios + opcionales del equipo)}
\label{tab:campos-json}
\end{longtable}

\section{Normalización y Tolerancia}
\label{sec:normalizacion}

El loader (\code{scenario/loader.py}) acepta y normaliza:

\begin{itemize}
\item \textbf{Bandos:} \code{"A"}/\code{"a"} $\to$ \code{Side.A}, \code{"B"}/\code{"b"} $\to$ \code{Side.B}
\item \textbf{Costos:} enteros o decimales positivos (\code{float})
\item \textbf{Terrenos no transitables:} \code{costo} opcional (se ignora si \code{transitable: false})
\item \textbf{Campos propios:} \code{juego.limite\_turnos}, \code{utilidad}, \code{prueba.*} (RNF-02)
\item \textbf{Variante dinámica:} \code{juego.variante}, \code{juego.turnos\_recuperacion},
  \code{juego.inmunidad\_robo}; si no se indica, se usa \code{base}
\end{itemize}

\section{Ejemplo Oficial (Apéndice A)}
\label{sec:ejemplo-apendice-a}

\begin{minted}{json}
{
  "version": "1.0",
  "mapa": { "filas": 4, "columnas": 4 },
  "tipos_terreno": {
    "camino":  {"costo": 1, "transitable": true},
    "pasto":   {"costo": 2, "transitable": true},
    "bosque":  {"costo": 4, "transitable": true},
    "pantano": {"costo": 7, "transitable": true},
    "muro":    {"transitable": false}
  },
  "terreno": [
    ["camino", "camino", "pasto",  "pasto"],
    ["camino", "muro",   "bosque", "pasto"],
    ["camino", "camino", "pantano","pasto"],
    ["pasto",  "pasto",  "camino", "camino"]
  ],
  "bases": { "A": {"fila": 0, "columna": 0}, "B": {"fila": 3, "columna": 3} },
  "recurso": {"fila": 2, "columna": 2},
  "unidades": [
    {"id": "A1", "bando": "A", "tipo": "estandar", "fila": 0, "columna": 1},
    {"id": "A2", "bando": "A", "tipo": "estandar", "fila": 1, "columna": 0},
    {"id": "B1", "bando": "B", "tipo": "estandar", "fila": 3, "columna": 2},
    {"id": "B2", "bando": "B", "tipo": "estandar", "fila": 2, "columna": 3}
  ],
  "turno": "A",
  "juego": { "portador_recurso": null, "profundidad_maxima_minimax": 4 },
  "prueba": { "modo": "busqueda", "unidad_inicio": "A1", "objetivo": {"fila": 3, "columna": 2} }
}
\end{minted}

\section{Validación Atómica (RF-34, RF-35)}
\label{sec:validacion-atomica}

\code{load\_scenario(path)} implementa una validación \textbf{atómica}:

\begin{enumerate}
\item Lee y parsea JSON completo
\item Ejecuta \textbf{todas} las validaciones, acumulando errores
\item Si hay \textbf{cualquier} error $\to$ lanza \code{ScenarioError} con \textbf{lista completa}
  de mensajes; \textbf{nada} se construye ni modifica
\item Si todo OK $\to$ construye \code{Scenario(board, initial\_state, test\_config)}
\end{enumerate}

\section{9 Validaciones Mínimas (RF-35) + 1 Extra}
\label{sec:9-validaciones}

\begin{longtable}{p{1cm}p{3cm}p{10cm}}
\toprule
\textbf{\#} & \textbf{Validación} & \textbf{Código / Detalle} \\
\midrule
\endfirsthead
\toprule
\textbf{\#} & \textbf{Validación} & \textbf{Código / Detalle} \\
\midrule
\endhead
\midrule
\multicolumn{3}{r}{\textit{Continúa}} \\
\endfoot
\bottomrule
\endlastfoot
1 & Filas/columnas positivas y consistentes con matriz & \code{\_read\_dimensions} + \code{\_read\_grid} \\
2 & Terrenos usados existen en \code{tipos\_terreno} & \code{\_read\_terrains} + \code{\_read\_grid} \\
3 & Transitables tienen costo positivo & \code{\_read\_terrains} \\
4 & Bases, recurso, unidades dentro del mapa & \code{\_read\_bases}, \code{\_read\_position}, \code{\_read\_units} \\
5 & Ninguna unidad en celda no transitable & \code{\_read\_units} \\
6 & IDs de unidad únicos & \code{\_read\_units} \\
7 & Cada unidad en bando A/a o B/b & \code{\_read\_side} \\
8 & Turno es A/a o B/b & \code{\_read\_side} \\
9 & \code{portador\_recurso} es id existente o null & \code{\_read\_carrier} \\
+ & \textbf{Extra:} Ninguna unidad comparte celda & \code{\_read\_units} (mapa \code{cells}) \\
\bottomrule
\caption{10 validaciones (9 RF-35 + 1 extra del equipo)}
\label{tab:10-validaciones}
\end{longtable}

\section{Campos de \code{prueba} (G10, Decisión 2026-10-17)}
\label{sec:campos-prueba}

El enunciado no fija nombres canónicos para algoritmo, heurística, $k$, profundidad, orden.
El equipo definió campos opcionales propios en \code{prueba} (RNF-02), CLI sobrepone JSON:

\begin{table}[H]
\centering
\caption{Campos opcionales en \code{prueba}}
\begin{tabular}{lll}
\toprule
\textbf{Campo} & \textbf{Valores} & \textbf{Descripción} \\
\midrule
\code{modo} & \code{busqueda}, \code{adversarial}, \code{partida} & Modo de ejecución \\
\code{algoritmo} & \code{BFS}, \code{DFS}, \code{UCS}, \code{A\_STAR}, \code{BEAM}, \code{MINIMAX}, \code{ALPHA\_BETA} & Algoritmo a ejecutar \\
\code{heuristica} & \code{H1}, \code{H2} (\code{MANHATTAN}, \code{OBSTACLE\_AWARE}) & Heurística para A*/Beam \\
\code{k} & int $\ge$1 & Ancho de haz (Beam) \\
\code{profundidad} / \code{depth} & int $\ge$1 & Profundidad Minimax/$\alpha$-$\beta$ \\
\code{orden\_acciones} & \code{default}, \code{reversed} & Orden de sucesores \\
\code{unidad\_inicio} & string & id unidad de inicio (búsqueda) \\
\code{objetivo} & \code{\{fila, columna\}} & Meta de búsqueda \\
\bottomrule
\end{tabular}
\end{table}
\label{tab:prueba-campos}


Aliases canónicos: \code{A*}/\code{A\_ESTRELLA} $\to$ \code{A\_STAR}, \code{H1}/\code{MANHATTAN},
\code{H2}/\code{OBSTACLE\_AWARE}. CLI (\code{--algorithm}, \code{--heuristic}, \code{--k},
\code{--depth}, \code{--order}) \textbf{sobrepone} al JSON (G10).

\section{Ejemplos de Escenarios del Equipo}
\label{sec:escenarios-equipo}

\begin{itemize}
\item \code{apendice\_a.json}: Ejemplo oficial 4$\times$4
\item \code{caso\_01\_bfs\_vs\_dfs.json} -- Caso 1
\item \code{caso\_02\_bfs\_vs\_ucs.json} -- Caso 2
\item \code{caso\_03\_obstaculo\_heuristica.json} -- Caso 3
\item \code{caso\_04\_comparacion\_heuristicas.json} -- Caso 4
\item \code{caso\_05\_beam\_k.json} -- Caso 5
\item \code{caso\_07\_decision\_adversarial.json} -- Casos 6--10 (mapa 5$\times$7)
\item \code{prueba\_15x15.json} -- Rendimiento 15$\times$15 (3 unidades/bando)
\item \code{prueba\_20x20.json} -- Prueba 20$\times$20
\item \code{escenario\_03x03} .. \code{escenario\_30x30} -- Escalabilidad
\end{itemize}
% 13_limitaciones.tex - Capítulo 13: Limitaciones
\chapter{Limitaciones y Trabajo Futuro}
\label{ch:limitaciones}

\section{Limitaciones Conocidas}
\label{sec:limitaciones-conocidas}

\begin{enumerate}
\item \textbf{Validación cruzada.} El equipo revisó los módulos, las pruebas y la
  documentación antes de la entrega.
  
\item \textbf{Pesos de utilidad con limitación conocida.} Con defaults 10/3/1/0.5:
  \begin{itemize}
  \item Recurso en celda propia de MAX (nadie porta, dist=0): $U = 0$
  \item MAX porta y está a 4 pasos de su base: $U = +3 - 1\cdot4 = -1$
  \item \textbf{Anomalía:} Recurso en celda propia ($U=0$) se valora \textbf{mejor} que
    llevarlo a más de 3 pasos de la base ($U < 0$).
  \item Profundidades pares aplanan valores a 0 (efecto horizonte, caso 10).
  \end{itemize}
  \textbf{Decisión (2026-10-17):} No cambiar defaults (invalidaría casos 7--10 ya validados).
  Se documenta como limitación; quien necesite otro comportamiento ajusta
  \code{recurso\_cargado $\ge$ 8} en el JSON de su escenario (\texttt{utilidad.recurso\_cargado}).

\item \textbf{Complejidad exponencial de Minimax.} Minimax puro crece exponencialmente:
  profundidad 6 en 15$\times$15 $\approx$ 15.7 s (supera meta 5 s). $\alpha$-$\beta$ mitiga
  (d=7 $\approx$ 0.86 s), pero depende del orden de acciones. En mapas grandes con muchas
  unidades, la poda es menos efectiva.

\item \textbf{Búsquedas de un agente ignoran al adversario.} Por diseño (RF-14), rivales se
  eliminan, turno no alterna, recurso ignorado. No modelan un plan de juego completo,
  solo navegación puntual.

\item \textbf{Solo dos acciones: Move y Wait.} No hay acciones de proteger, bloquear, atacar
  u otras (decisión 2026-10-07: complejidad innecesaria). El placaje es consecuencia de
  movimiento, no una acción explícita.

\item \textbf{Estados explorados no visualizados en GUI.} RF-41 es opcional; la GUI muestra
  camino/acción sugerida, no frontera/estados explorados.

\item \textbf{Preguntas al profesor sin respuesta oficial.} Las 8 preguntas de la Sección 19
  del enunciado (vecindad, semántica de costo, campos \code{prueba}, mayúsculas, límites,
  ocupación, decimales, formato entrega) no tienen respuesta oficial; las decisiones del
  equipo están documentadas en \texttt{decisions.md}.
\end{enumerate}

\section{Trabajo Futuro}
\label{sec:trabajo-futuro}

\begin{enumerate}
\item \textbf{Revisar pesos de utilidad:} Ajustar \code{recurso\_cargado} $\ge$ 8 o añadir
  término de costo de terreno para evitar anomalía del recurso en celda propia.
\item \textbf{Ordenamiento de acciones en $\alpha$-$\beta$:} Implementar ordenación por
  heurística (mejores primero) para maximizar poda.
\item \textbf{Transposition table en Minimax:} Memoizar estados evaluados para evitar
  recálculos en árboles con transposiciones.
\item \textbf{Visualización de frontera/expandidos en GUI:} Cumplir RF-41 opcional.
\item \textbf{Acciones adversariales adicionales:} Proteger zona, bloquear posición,
  interceptar explícito (requiere ampliar \code{Action} y \code{legal\_actions}).
\item \textbf{Tests de estrés y fuzzing JSON:} Generar escenarios aleatorios válidos para
  validar robustez del loader.
\item \textbf{CI/CD:} GitHub Actions para \texttt{pytest}, \texttt{ruff}, \texttt{mypy} en
  cada push.
\item \textbf{Documentación en inglés:} Traducir .md y docstrings (código ya en inglés).
\end{enumerate}
% 14_uso_ia.tex - Capítulo 14: Uso de IA Generativa
\chapter{Reporte de Uso de IA Generativa}
\label{ch:ia}

Este capítulo presenta de forma transparente el apoyo recibido de una herramienta de
inteligencia artificial en la redacción, organización, programación y revisión del trabajo.
La herramienta fue un apoyo para el equipo y no reemplazó el análisis, las decisiones,
la implementación ni la validación final del proyecto.

\section{Herramienta}
\label{sec:herramienta-ia}

\textbf{Herramienta de apoyo:} \textbf{Claude Code} (modelo \textbf{Claude Sonnet 5.5}),
utilizada por el equipo como apoyo para consultas técnicas, redacción, organización y
revisión del material. La herramienta no fue el desarrollador principal ni tomó decisiones
autónomas sobre el proyecto.

\section{Propósito y Áreas de Uso}
\label{sec:proposito-ia}

\begin{longtable}{p{3cm}p{3cm}p{8cm}}
\toprule
\textbf{Área} & \textbf{Qué hizo la IA} & \textbf{Qué hizo el equipo} \\
\midrule
\endfirsthead
\toprule
\textbf{Área} & \textbf{Qué hizo la IA} & \textbf{Qué hizo el equipo} \\
\midrule
\endhead
\midrule
\multicolumn{3}{r}{\textit{Continúa}} \\
\endfoot
\bottomrule
\endlastfoot
Contexto y proceso & Sugirió una organización de archivos y fases de trabajo & El equipo evaluó, adaptó y registró las decisiones en \texttt{decisions.md} \\
Decisiones & Presentó alternativas técnicas para discusión & El equipo eligió las reglas, el alcance y las convenciones del sistema \\
\code{domain/}, \code{scenario/} & Apoyó la redacción de borradores de código a partir de las reglas definidas & El equipo definió las reglas, revisó la implementación y exigió validaciones adicionales \\
BFS, DFS, UCS, A*, Beam, verifier & Apoyó la elaboración de código de referencia y pruebas & El equipo definió las convenciones y derivó los valores esperados antes de ejecutar \\
Heurísticas & Apoyó la exploración de alternativas y la redacción de ambas heurísticas & El equipo verificó admisibilidad, consistencia y dominancia h2$\ge$h1 \\
Utilidad, Minimax, $\alpha$-$\beta$ & Apoyó borradores de código y fórmulas para discusión & El equipo aprobó, corrigió y validó la utilidad y los algoritmos \\
CLI, runner, sweeps, export & Apoyó ejemplos y revisión de interfaces & El equipo decidió las opciones y los formatos CSV/JSON \\
GUI & Apoyó borradores de \code{gui/session.py} y \code{gui/app.py} & El equipo definió las funciones, revisó el comportamiento y validó la interfaz \\
Escenarios y experimentos & Apoyó la generación de mapas y la organización de resultados & El equipo escogió los casos y calculó los valores esperados \\
Documentación & Apoyó la redacción, corrección de estilo y organización de este informe y del manual & El equipo revisó, corrigió y asumió la revisión final del contenido \\
\bottomrule
\caption{Uso de IA por área}
\label{tab:uso-ia}
\end{longtable}

\section{Intervenciones Humanas y Correcciones}
\label{sec:intervenciones-humanas}

Las siguientes correcciones humanas se realizaron sobre los borradores de la IA:

\begin{enumerate}
\item \textbf{Signo del término rival (utilidad):} La decisión textual daba un único término
  negativo rival; la implementación hace que el signo dependa de quién porta el recurso
  (un rival cerca del \textit{portador de MAX} baja $U$; cerca del portador de MIN no suma).
\item \textbf{Rivales como obstáculos en búsqueda 1-agente:} El primer modelo trataba rivales
  como obstáculos fijos, lo que dejaba el Apéndice A sin solución (B1 sobre la meta). El
  equipo cambió a: \textbf{quitar rivales}, compañeros quietos como obstáculos (Sección 6.1).
\item \textbf{Inconsistencia de documentación:} Se revisaron los textos que describían el
  origen de las propuestas para distinguir claramente entre sugerencias de la herramienta
  y decisiones adoptadas por el equipo.
\item \textbf{Validación de casos 1--10:} Valores esperados derivados \textbf{a mano del mapa
  antes de ejecutar}; tests asercionan esos valores. La IA no inventó los resultados.
\item \textbf{Revisión humana:} El equipo revisó y validó los resultados, el código y la
  documentación generados con apoyo de herramientas de IA.
\end{enumerate}

La IA se utilizó como apoyo técnico y editorial; las decisiones de diseño, la
implementación, la verificación y la responsabilidad final del trabajo corresponden al
equipo.
% 15_reproducibilidad.tex - Capítulo 15: Reproducibilidad y Versiones
\chapter{Reproducibilidad, Versiones y Entrega}
\label{ch:reproducibilidad}

\section{Instalación y Ejecución}
\label{sec:instalacion}

\begin{minted}{bash}
# 1. Clonar repositorio
git clone https://github.com/df1903/tactical-grid-ai
cd tactical-grid-ai

# 2. Instalar dependencias (uv provisiona Python 3.14)
uv sync

# 3. Verificar instalación
uv run pytest              # 281 tests deben pasar
uv run ruff check .        # lint
uv run mypy src            # tipos strict
\end{minted}

\section{Versiones Fijas (RE-02, RE-03, RE-04)}
\label{sec:versiones}

Todas las dependencias fijas en \code{uv.lock}:

\begin{table}[H]
\centering
\caption{Versiones fijas de dependencias}
\begin{tabular}{ll}
\toprule
\textbf{Herramienta} & \textbf{Versión} \\
\midrule
Python & 3.14 (provisionado por \code{uv}) \\
\texttt{uv} & 0.x (gestor de paquetes) \\
\texttt{pytest} & 9.1.1 \\
\texttt{ruff} & 0.16.10 (reglas E, F, I, UP, B) \\
\texttt{mypy} & 2.4.0 (--strict) \\
\texttt{minted} & (LaTeX) \\
\texttt{pgfgantt}, \texttt{pgf-pie} & (LaTeX diagrams) \\
\bottomrule
\end{tabular}
\end{table}
\label{tab:versiones}


La entrega se acompaña del repositorio Git con el código fuente, los escenarios,
los resultados y la documentación utilizados para generar estos documentos.

\section{Estados Iniciales Declarados}
\label{sec:estados-iniciales}

Todos los estados iniciales están declarados en archivos JSON de \code{scenarios/}:

\begin{itemize}
\item \code{apendice\_a.json} -- Ejemplo oficial 4$\times$4
\item \code{caso\_01\_bfs\_vs\_dfs.json} -- Caso 1
\item \code{caso\_02\_bfs\_vs\_ucs.json} -- Caso 2
\item \code{caso\_03\_obstaculo\_heuristica.json} -- Caso 3
\item \code{caso\_04\_comparacion\_heuristicas.json} -- Caso 4
\item \code{caso\_05\_beam\_k.json} -- Caso 5
\item \code{caso\_07\_decision\_adversarial.json} -- Casos 6--10 (mapa 5$\times$7)
\item \code{prueba\_15x15.json} -- Rendimiento 15$\times$15 (3 unidades/bando)
\item \code{prueba\_20x20.json} -- Mapa 20$\times$20
\item \code{escenario\_03x03} .. \code{escenario\_30x30} -- Escalabilidad
\end{itemize}

\section{Resultados Esperados vs. Obtenidos}
\label{sec:esperado-vs-obtenido}

Todos los 10 casos mínimos tienen sus valores esperados (derivados a mano) y obtenidos
(ejecución real) documentados en \texttt{cases.md} y \texttt{docs/results/}.
Los tests \code{test\_cases\_f2.py}, \code{test\_cases\_f3.py}, \code{test\_case\_6.py},
\code{test\_cases\_f4.py} asercionan \textbf{expected == obtained}.

Los archivos CSV/JSON brutos están en \texttt{docs/results/}:

\begin{itemize}
\item \code{caso\_01\_bfs\_vs\_dfs.csv/json}
\item \code{caso\_02\_bfs\_vs\_ucs.csv/json}
\item \code{caso\_03\_obstaculo\_heuristica.csv/json}
\item \code{caso\_04\_comparacion\_heuristicas.csv/json}
\item \code{caso\_05\_beam\_k.csv/json}
\item \code{casos\_07\_10\_adversarial.csv/json}
\item \code{rendimiento\_15x15.csv}, \code{rendimiento\_20x20.csv}
\item \code{apendice\_a.csv/json}
\end{itemize}

\section{Comandos de Reproducción (Evidencia RE-03, RE-04)}
\label{sec:comandos-reproduccion}

Todos los comandos exactos para regenerar los resultados están en
\texttt{user-manual.md} §6 y se resumen aquí:

\begin{minted}{bash}
R=docs/results; U="uv run python -m tactical_grid"

# Casos 1-2 (búsquedas no informadas / costos)
$U scenarios/caso_01_bfs_vs_dfs.json --algorithm BFS,DFS,UCS --verify --csv $R/caso_01_bfs_vs_dfs.csv --json $R/caso_01_bfs_vs_dfs.json
$U scenarios/caso_02_bfs_vs_ucs.json --algorithm BFS,DFS,UCS --verify --csv $R/caso_02_bfs_vs_ucs.csv --json $R/caso_02_bfs_vs_ucs.json

# Casos 3-5 (informada)
$U scenarios/caso_03_obstaculo_heuristica.json --algorithm UCS,A_STAR --heuristic H1,H2 --verify --csv $R/caso_03_obstaculo_heuristica.csv --json $R/caso_03_obstaculo_heuristica.json
$U scenarios/caso_04_comparacion_heuristicas.json --algorithm UCS,A_STAR --heuristic H1,H2 --verify --csv $R/caso_04_comparacion_heuristicas.csv --json $R/caso_04_comparacion_heuristicas.json
$U scenarios/caso_05_beam_k.json --algorithm UCS,BEAM --heuristic H1 --k 1,2,4,8 --verify --csv $R/caso_05_beam_k.csv --json $R/caso_05_beam_k.json

# Casos 6-10 (adversarial)
$U scenarios/caso_07_decision_adversarial.json --algorithm MINIMAX --depth 1,2,3 --verify --csv $R/casos_07_10_adversarial.csv --json $R/casos_07_10_adversarial.json
$U scenarios/caso_07_decision_adversarial.json --algorithm MINIMAX,ALPHA_BETA --depth 4 --verify --csv $R/casos_07_10_adversarial.csv --json $R/casos_07_10_adversarial.json
$U scenarios/caso_07_decision_adversarial.json --algorithm ALPHA_BETA --depth 3 --order default,reversed --verify --csv $R/casos_07_10_adversarial.csv --json $R/casos_07_10_adversarial.json
$U scenarios/caso_07_decision_adversarial.json --algorithm ALPHA_BETA --depth 1,2,3,4,5 --verify --csv $R/casos_07_10_adversarial.csv --json $R/casos_07_10_adversarial.json

# Rendimiento
$U scenarios/prueba_20x20.json --algorithm ALPHA_BETA --depth 2,3,4,5,6 --csv $R/rendimiento_20x20.csv
$U scenarios/prueba_15x15.json --algorithm ALPHA_BETA --depth 5,6,7 --csv $R/rendimiento_15x15.csv
$U scenarios/prueba_15x15.json --algorithm MINIMAX --depth 4,5,6 --csv $R/rendimiento_15x15_minimax.csv

# Sweep completo Apéndice A
$U scenarios/apendice_a.json --algorithm BFS,DFS,UCS,A_STAR,BEAM --heuristic H1,H2 --k 1,2,4,8 --verify --csv $R/apendice_a.csv
\end{minted}

\section{Identificación de la Entrega (RE-02)}
\label{sec:tag-entrega}

La entrega puede identificarse y reproducirse desde el repositorio del proyecto:

\begin{minted}{bash}
git describe --tags --always
\end{minted}

El comando anterior permite consultar el identificador exacto del estado entregado.

\section{Checklist Final de Entrega}
\label{sec:checklist-entrega}

\begin{longtable}{ll}
\toprule
\textbf{Ítem} & \textbf{Estado} \\
\midrule
\endfirsthead
\toprule
\textbf{Ítem} & \textbf{Estado} \\
\midrule
\endhead
\midrule
\multicolumn{2}{r}{\textit{Continúa}} \\
\endfoot
\bottomrule
\endlastfoot
Repositorio Git con aportes visibles por integrante & \checkmark \\
README.md con nombres de integrantes & \checkmark \\
Código fuente completo + instrucciones instalación/ejecución & \checkmark \\
Escenarios JSON usados + casos mínimos de prueba & \checkmark \\
Manual de usuario completo & \checkmark \\
Documento técnico completo (incluye limitaciones, reporte IA) & \checkmark \\
Escenarios JSON + 10 casos con esperado vs. obtenido & \checkmark \\
Repositorio con aportes visibles e identificador reproducible & \checkmark \\
Uso de IA declarado (herramienta, propósito, intervención, validación) & \checkmark \\
Cada integrante puede explicar cualquier parte & \checkmark \\
\bottomrule
\end{longtable}

% 16_conclusiones.tex - Capítulo 16: Conclusiones
\chapter{Conclusiones}
\label{ch:conclusiones}

\section{Resumen de Logros}
\label{sec:logros}

El proyecto \tacticalgrid\ cumple \textbf{todos los requisitos obligatorios} del enunciado del
Trabajo de Sistemas Inteligentes I (UCALDAS, 2026-2):

\begin{itemize}
\item \textbf{Modelado y estado (RF-10--13, caso 6):} Estado compuesto inmutable, reglas
  deterministas, identificación unívoca, justificación de suficiencia.
\item \textbf{BFS, DFS, UCS (RF-14--19, casos 1--2):} 6 métricas, convenciones G7/G8,
  casos 1--2 con esperado vs. obtenido.
\item \textbf{A*, Beam Search (RF-20--24, casos 3--5):} 2 heurísticas (Manhattan escalada,
  obstacle-aware BFS inverso), justificación RF-22 completa, Beam k=1,2,4,8.
\item \textbf{Minimax, $\alpha$-$\beta$ (RF-25--33, casos 7--10):} Utilidad con pesos JSON,
  7 métricas RF-30, misma decisión (RF-31), análisis orden/poda (RF-32), check 6 verif.
\item \textbf{Experimentación (RF-16, 29, 30, 38, 39, RE-03/04):} 10 casos con esperado vs
  obtenido, sweep CLI, export CSV/JSON, rendimiento medido.
\item \textbf{Validación (RF-34, 35, 42):} Carga atómica 9+1 validaciones, 6 checks verif.
\item \textbf{Calidad y docs (RNF-05, RD-01, 9.2):} 5 componentes, docstrings completos,
  281 tests, ruff/mypy strict, reporte IA.
\item \textbf{GUI (RF-40):} tkinter, colores por costo JSON, auto IA vs IA, carga JSON atómico.
\item \textbf{CLI + Sweep (G10):} Opciones listas, sweep combinatorio, export CSV/JSON.
\end{itemize}

\section{Aportes Destacados}
\label{sec:aportes}

\begin{enumerate}
\item \textbf{Heurística obstacle-aware (h2):} BFS inverso desde la meta que ve muros y base
  rival. Domina a Manhattan (h2 $\ge$ h1), admisible y consistente. Reduce expandidos
  2$\times$ en caso 3, 1.6$\times$ en caso 4.
\item \textbf{Placaje como acción explícita:} Reemplaza interceptación automática. Acción
  \code{Plack} con costo = terreno destino; recurso cae en celda del portador (sin transferencia).
  Elimina desempates y simplifica reglas.
\item \textbf{Utilidad con objetivo dinámico y término rival condicional:} Objetivo cambia según
  quién porta; término rival solo cuando MAX porta. Pesos 100\% configurables en JSON (RNF-03).
\item \textbf{Abstracción 1-agente limpia:} \code{make\_problem} quita rivales, fija turno,
  quita recurso; \code{\_normalize} limpia side-effects de \code{apply}.
\item \textbf{Sweep CLI + export CSV/JSON:} Listas comma-separated, producto cartesiano de
  combos, export plano (CSV) y completo (JSON con \code{valores\_raiz}).
\item \textbf{GUI tkinter única:} 700 líneas, colores por costo JSON, auto IA vs IA, deshacer,
  reiniciar, carga atómica, ayuda contextual.
\end{enumerate}

\section{Lecciones Aprendidas}
\label{sec:lecciones}

\begin{enumerate}
\item \textbf{Definir convenciones antes de codificar:} G7, G8, G9 (tie-break, meta al expandir,
  Beam criterion) decidieron la corrección de métricas.
\item \textbf{Escalar heurísticas por costo mínimo del escenario:} Garantiza admisibilidad
  ante mapas externos con costos $< 1$ (enunciado §18.4).
\item \textbf{Placaje como acción, no side-effect:} Simplifica \code{apply}, elimina
  desempates, hace reglas más explicables.
\item \textbf{Utilidad con pesos en JSON:} Cumple RNF-03 sin hardcodeo; permite ajustar
  comportamiento sin tocar código.
\item \textbf{Validar a mano antes de ejecutar:} Valores esperados de casos 1--10 derivados
  a mano $\to$ tests asercionan \texttt{expected == obtained}.
\item \textbf{Documentar decisiones en \texttt{decisions.md}:} 16 decisiones registradas con
  fecha, contexto, alternativas descartadas; trazabilidad total.
\end{enumerate}

\section{Trabajo Futuro}
\label{sec:trabajo-futuro-final}

Ver Cap. \ref{ch:limitaciones} §\ref{sec:trabajo-futuro}. Las líneas de trabajo
identificadas para futuras ampliaciones son:
\begin{enumerate}
\item Evaluar ajustes de \code{recurso\_cargado} en la función de utilidad.
\item Ampliar la interfaz con más visualizaciones y análisis comparativos.
\item Incorporar nuevos escenarios para evaluar escalabilidad.
\end{enumerate}

\section{Agradecimientos}
\label{sec:agradecimientos}

A la Universidad de Caldas, al curso Sistemas Inteligentes I, y al profesor por la guía.
Al modelo \textbf{Claude Sonnet 5.5} (Claude Code) por asistir en redacción de código,
documentación y decisiones; las intervenciones humanas y responsabilidad final son del
equipo (Jorge Iván Garcia Torres, Daniel Felipe Franco Rincón).

\vspace{2cm}
\begin{flushright}
\textbf{Manizales, 21 de octubre de 2026}\\
\vspace{1cm}
Jorge Iván Garcia Torres \hspace{3cm} Daniel Felipe Franco Rincón
\end{flushright}

% Bibliografía
\clearpage
\printbibliography[heading=bibintoc,title={Referencias Bibliográficas}]

% Apéndices
\appendix
\chapter{Diagramas Mermaid $\to$ TikZ}
\label{ap:diagramas}
Los diagramas del sistema se presentan como gráficos vectoriales, ajustados
proporcionalmente al área útil de la página para evitar recortes.

\diagramfigure{arquitectura}{0.96}{Arquitectura de componentes y flujo de datos.}
\diagramfigure{ciclo_turno}{0.78}{Ciclo de turno y condiciones de victoria.}
\diagramfigure{flujo_apply}{0.86}{Flujo de aplicación de una acción.}
\diagramfigure{flujo_apply_intercept}{0.86}{Flujo de aplicación con interceptación.}
\diagramfigure{journey_progreso}{0.96}{Cronograma de progreso del proyecto.}

\chapter{Código Fuente Completo (Listings)}
\label{ap:codigo}
\section{Módulo \code{domain/rules.py}}
\inputminted{python}{../../src/tactical_grid/domain/rules.py}

\section{Módulo \code{algorithms/informed.py}}
\inputminted{python}{../../src/tactical_grid/algorithms/informed.py}

\section{Módulo \code{algorithms/adversarial.py}}
\inputminted{python}{../../src/tactical_grid/algorithms/adversarial.py}

\section{Módulo \code{algorithms/utility.py}}
\inputminted{python}{../../src/tactical_grid/algorithms/utility.py}

\section{Módulo \code{scenario/loader.py}}
\inputminted{python}{../../src/tactical_grid/scenario/loader.py}

\section{Módulo \code{experiments/verifier.py}}
\inputminted{python}{../../src/tactical_grid/experiments/verifier.py}

\chapter{Resultados Completos (CSV/JSON)}
\label{ap:resultados}
Los archivos CSV y JSON completos de todos los experimentos se encuentran en \texttt{docs/results/}.

\end{document}